%% relatorio_1610.tex
%% Relatorio técnico de respuesta a las Orientações de 03/10/2026 (Prof. Hugo E. Hernández-Figueroa)
%% Experimentos A-F: canal monoestático permanente + canal biestático complementario,
%% fusión coherente normalizada, pseudo-Brewster como supresión de superficie.
%% Código: simulations/analise_1610 (README.md). Datos: simulations/io/analise_1610/expX/*.csv
%% Compilar: pdflatex relatorio_1610.tex (dos pasadas)

\documentclass[11pt,a4paper]{report}

\usepackage[utf8]{inputenc}
\usepackage[T1]{fontenc}
\usepackage[spanish,es-tabla,es-noshorthands]{babel}
\usepackage{lmodern}
\usepackage{geometry}
\geometry{left=2.5cm, right=2.3cm, top=2.6cm, bottom=2.6cm}
\usepackage{amsmath,amssymb}
\usepackage{bm}
\usepackage{booktabs}
\usepackage{array}
\usepackage{longtable}
\usepackage{multirow}
\usepackage{enumitem}
\usepackage{xcolor}
\usepackage{colortbl}
\definecolor{mejor}{HTML}{CDEFD9}
\definecolor{regular}{HTML}{FFF1C2}
\definecolor{peor}{HTML}{F8D3D0}
% conclusión de una frase al final de cada experimento
\newcommand{\conclexp}[2]{\par\medskip\noindent\colorbox{mejor!55}{\parbox{\dimexpr\linewidth-2\fboxsep\relax}{%
\textbf{Conclusión del Exp.~#1.} #2}}\par\medskip}
\usepackage{graphicx}
\usepackage{caption}
\usepackage{subcaption}
\usepackage{float}
\usepackage[colorlinks=true,linkcolor=blue!50!black,citecolor=blue!50!black,urlcolor=blue!50!black]{hyperref}
\usepackage{titlesec}

\graphicspath{{figures/relatorio_1610/}}
\captionsetup{font=small, labelfont=bf}
\titleformat{\chapter}[display]{\normalfont\huge\bfseries}{\chaptertitlename\ \thechapter}{12pt}{\Huge}
\titlespacing*{\chapter}{0pt}{-20pt}{25pt}
\titleformat{\section}{\normalfont\Large\bfseries}{\thesection}{1em}{}
\titleformat{\subsection}{\normalfont\large\bfseries}{\thesubsection}{1em}{}
\setlength{\parindent}{0pt}
\setlength{\emergencystretch}{3em}
\setlength{\parskip}{0.5\baselineskip}

\newcommand{\dB}{\,\mathrm{dB}}
\newcommand{\vect}[1]{\bm{#1}}
\newcolumntype{L}[1]{>{\raggedright\arraybackslash}p{#1}}
\newcolumntype{C}[1]{>{\centering\arraybackslash}p{#1}}

\title{\vspace{-1cm}\textbf{SAR biestático subsuperficial:\\
canal biestático complementario al monoestático,\\
fusión coherente y supresión de superficie por pseudo-Brewster}\\[10pt]
\large Relatorio de respuesta a las Orientações del 03/10/2026}
\author{David Atencia}
\date{16 de octubre de 2026}

\begin{document}
\sloppy
\maketitle

\begin{abstract}
\noindent
Este relatorio responde a las orientaciones del 03/10/2026: el Tx es un
\emph{transceiver} que siempre adquiere el canal monoestático, y la pregunta es cuándo,
dónde y por qué un canal biestático simultáneo mejora esa imagen. Se construyó un motor
de simulación coherente (ecos complejos del mismo pulso, refracción de Fermat/Snell,
suelo complejo, Óptica Física biestática, superficie especular y difusa SPM,
back-projection 3D), validado contra el pipeline MATLAB existente, contra la Tabla~2 del
relatorio del 28/09 y contra soluciones analíticas. Con él se ejecutaron los
experimentos A--F y se calcularon todas las métricas pedidas en CSV.
La fusión se normaliza de forma explícita: una fusión con el mismo número de muestras
que la mono da exactamente la misma imagen, de modo que toda mejora de forma es
geométrica.
Resultados principales. (1) Ref1 (diametral) no enfoca en horizontal y su fusión empeora
la mono. (2) El canal bi a la misma altura contrae el anillo horizontal del dominio $k$
por $|\cos(\Delta\phi/2)|$. Con $\Delta\phi\approx90^{\circ}$--$120^{\circ}$, o con un barrido
$60^{\circ}\to150^{\circ}$, la fusión rellena el disco: ISLR $-5$ a $-8$\,dB, PSLR hasta $-5$\,dB y
contraste $+4$ a $+5.5$\,dB, a costa de un 16--30\,\% más de anchura $-3$\,dB. No mejora la
resolución vertical, que fija la banda. (3) Con objetos reales, un Rx TM co-polar pierde
$\approx20$\,dB a 90° por desajuste geométrico de polarización. Con un Rx de doble
polarización alineado, un Monte Carlo SO(3) por adquisición muestra una mejora de
contraste de $+1.8$ a $+2.7$\,dB en el 90--98\,\% de las orientaciones. (4) La fusión
coherente exige calibrar una fase entre canales: 1\,ns de offset sin calibrar la anula,
y una vez calibrada se recupera exactamente. (5) El pseudo-Brewster en suelo con pérdidas
($\varepsilon_r=3.99-j0.40$) es un mínimo ancho pero somero ($-34.5$\,dB) que sólo reduce
la especular coherente de Ref1: el SCR mejora 14--16\,dB pero queda en $-41$\,dB. En el
resto de las geometrías el clutter es difuso y no depende de Brewster. (6) Lo que sí suprime
el clutter es un nulo \emph{biestático} del retorno difuso VV: en TM, con
$\Delta\phi\approx135^{\circ}$--$140^{\circ}$, el canal bi \emph{aislado} mejora el SCR de una esfera y
de un cubo con orientación aleatoria en 8 a 20\,dB respecto de la mono, mientras que la
fusión de igual peso pierde la mayor parte de esa ganancia. El nulo HH en TE a 90°, que con
el punto escalar daba el mejor SCR, desaparece con blancos reales porque también desacopla
su eco. En el suelo sin
pérdidas el mínimo desaparece con 0.5° de error. La restricción cuesta poco en imagen
(${\pm}12$\,\%), pero no compensa. Se proponen dos modos de uso del Rx: para imagen, misma altura,
radio $0.85\rho_{Tx}$ y separación barrida de 60° a 150° con fusión coherente; para detección
limitada por clutter, separación fija de $\approx140^{\circ}$ en TM con el canal bi aislado o una
fusión ponderada. Se listan explícitamente los resultados negativos.
\end{abstract}


\tableofcontents

% ══════════════════════════════════════════════════════════════════════════
\chapter*{Glosario}
\addcontentsline{toc}{chapter}{Glosario}
\label{cap:glosario}

\section*{Siglas}

\begin{longtable}{@{}L{2.4cm}L{12.0cm}@{}}
\toprule
\textbf{Sigla} & \textbf{Significado} \\
\midrule
\endfirsthead
\toprule
\textbf{Sigla} & \textbf{Significado} \\
\midrule
\endhead
AIEM, IEM & \emph{(Advanced) Integral Equation Model}: modelos de dispersión de superficies rugosas
            más generales que el SPM; incluyen términos de orden superior. \\
BP & \emph{Back-projection}: reconstrucción que suma, en cada vóxel, los ecos corregidos por el
     retardo de ese vóxel. \\
CSV & Archivo de texto con columnas separadas por comas; contiene todas las métricas de cada
      experimento. \\
FWHM & \emph{Full Width at Half Maximum}: ancho a media potencia. Equivale a $\delta_{x,y,z}$
       (ancho a $-3$\,dB); los CSV conservan el nombre \texttt{fwhm\_*}. \\
GNSS, RTK & Posicionamiento por satélite; RTK (\emph{Real-Time Kinematic}) es la variante con
            precisión centimétrica. \\
GPSDO, OCXO & Osciladores de alta estabilidad (disciplinado por GPS; de cristal con horno) para
              sincronizar los relojes del Tx y del Rx. \\
HH, VV & Polarización horizontal (TE) o vertical (TM) en transmisión y recepción; se usa para el
         retorno de la superficie. \\
ISLR & \emph{Integrated Sidelobe Ratio}: energía fuera del lóbulo principal dividida por la
       energía dentro. Menor es mejor. \\
MIP & \emph{Maximum Intensity Projection}: proyección de una imagen 3D en 2D tomando el máximo de
      $|I|$ a lo largo de un eje. ``MIP en $z$'' da una vista en planta. \\
PEC & \emph{Perfect Electric Conductor}: conductor perfecto; material de los blancos simulados. \\
PO & \emph{Physical Optics} (Óptica Física): aproximación para calcular el campo dispersado por
     un objeto a partir de las corrientes en sus caras iluminadas. \\
PSCR & Pico del blanco dividido por el pico de la superficie en toda la rejilla (Exp.~C2). \\
PSF & \emph{Point Spread Function}: imagen de un único punto; describe la resolución y los
      lóbulos del sistema. \\
PSLR & \emph{Peak Sidelobe Ratio}: pico del mayor lóbulo secundario respecto del pico principal.
       Menor es mejor. \\
RCS & \emph{Radar Cross Section}: sección eficaz radar del blanco, en m$^2$ o dBsm. \\
Rx, Tx & Receptor y transmisor. El Tx es un \emph{transceiver}: también recibe el canal mono. \\
SAR & \emph{Synthetic Aperture Radar}: radar de apertura sintética. \\
SCR & \emph{Signal-to-Clutter Ratio}: potencia del blanco dividida por la de la superficie en el
      vóxel del blanco. Mayor es mejor. \\
SNR & \emph{Signal-to-Noise Ratio}: potencia del blanco dividida por la del ruido térmico. \\
SO(3) & Grupo de rotaciones en 3D. ``Monte Carlo SO(3)'': orientaciones del objeto sorteadas
        uniformemente. \\
SPM & \emph{Small Perturbation Method}: modelo de dispersión de primer orden para superficies
      poco rugosas ($k_0s\ll1$). \\
TM, TE & Polarización con el campo eléctrico en el plano de incidencia (TM, ``vertical'') o
         perpendicular a él (TE, ``horizontal''). Brewster sólo afecta a TM. \\
UAV & \emph{Unmanned Aerial Vehicle}: dron que lleva el Tx o el Rx. \\
\bottomrule
\end{longtable}

\section*{Configuraciones y nombres usados en el relatorio}

\begin{longtable}{@{}L{2.8cm}L{11.6cm}@{}}
\toprule
\textbf{Término} & \textbf{Significado} \\
\midrule
\endfirsthead
\toprule
\textbf{Término} & \textbf{Significado} \\
\midrule
\endhead
Mono, Ref2 & Canal monoestático: el Tx recibe su propio eco. Es la referencia permanente. \\
Bi & Canal biestático: el Rx remoto recibe el eco del mismo pulso del Tx. \\
Ref1 & Canal bi diametral: Rx a la misma altura y radio, en el lado opuesto ($\Delta\phi=180^{\circ}$). \\
fus (2N) & Fusión coherente $\tfrac12(I_{mono}+I_{bi})$ con los dos canales completos. \\
fusN (N) & Fusión con el mismo número de muestras que la mono (canales alternados por pulso).
           Aísla la ganancia geométrica de la ganancia por más muestras. \\
Fusión incoherente & Suma de intensidades, $\tfrac12(|I_{mono}|^2+|I_{bi}|^2)$; no necesita fase. \\
Hélice base & Trayectoria óptima del Tx del trabajo anterior: 2 vueltas, incidencia de 46° a 60°. \\
F1, F2 & Familias de trayectoria del Exp.~C: F1 es la hélice girada hacia una incidencia central;
         F2 es una órbita de incidencia constante. \\
S0--S3 & Suelos simulados: S0 sin pérdidas ($\varepsilon_r=4$), S1 de referencia
         ($3.99-j0.40$), S2 ($6.5-j0.9$) y S3 ($12-j2.5$), cada vez más húmedos. \\
P0, P1, P2 & Conjuntos de blancos del Exp.~D: objetos simples, objetos alargados o inclinados, y
             escena con tres objetos. \\
Rx alineado & Rx de doble polarización que combina sus canales para recibir la componente del
              campo del Tx; evita el desajuste de polarización a $\Delta\phi\approx90^{\circ}$. \\
Pseudo-Brewster & Incidencia con mínimo (no nulo) de la reflectancia TM en un suelo con pérdidas. \\
Nulo VV biestático & Separación $\Delta\phi^\ast\approx135^{\circ}$--$140^{\circ}$ donde se anula el retorno
                     difuso VV de primer orden. \\
Especular & Reflexión coherente de la superficie en el punto especular del par Tx--Rx; en el
            canal mono, el reflejo de nadir. \\
Difusa & Retorno incoherente de la rugosidad de la superficie (modelo SPM). \\
Dominio $k$ & Espacio de los vectores de onda $\vect K$ medidos; su ocupación fija la resolución
              y los lóbulos de la imagen. \\
Rango efectivo & Número de valores singulares significativos de la matriz de adquisición;
                 mide cuánta información independiente aporta la geometría. \\
Ponderación Hann & Ventana aplicada al espectro en frecuencia para bajar los lóbulos en rango,
                   a costa de un lóbulo principal más ancho. \\
Roll, atitud & Giro de la antena alrededor de la línea de vista; mezcla las polarizaciones TM y TE. \\
P$_{mej}$ & Probabilidad de mejora: fracción de adquisiciones del Monte Carlo en que la fusión
            supera a la mono. \\
\bottomrule
\end{longtable}

\section*{Símbolos}

\begin{longtable}{@{}L{2.8cm}L{11.6cm}@{}}
\toprule
\textbf{Símbolo} & \textbf{Significado} \\
\midrule
\endfirsthead
\toprule
\textbf{Símbolo} & \textbf{Significado} \\
\midrule
\endhead
$\Delta\phi$ & Separación acimutal entre Tx y Rx, vista desde el centro de la escena. \\
$\kappa$ & Razón entre el radio de la órbita del Rx y la del Tx. \\
$\rho$, $z$ & Radio horizontal y altura de una plataforma. \\
$\theta$, $\theta_c$ & Incidencia en la superficie; $\theta_c$ es la incidencia central de una
                       trayectoria del Exp.~C. \\
$\theta_{pB}$ & Incidencia de pseudo-Brewster del suelo. \\
$\theta_{esp}$ & Incidencia en el punto especular del par Tx--Rx. \\
$\delta_x,\delta_y,\delta_z$ & Resoluciones: ancho a $-3$\,dB de la PSF en cada eje. \\
$V_{-3}$ & Volumen del lóbulo principal a $-3$\,dB, en litros. \\
$\varepsilon_r$ & Permitividad relativa compleja del suelo, $\varepsilon'-j\varepsilon''$. \\
$n$ & Índice de refracción del suelo, $\sqrt{\varepsilon_r}$. \\
$s$, $\ell$ & Rugosidad de la superficie: desviación típica de la altura y longitud de correlación. \\
$\sigma^0$ & Coeficiente de dispersión de la superficie por unidad de área. \\
$r_{TM}$, $r_{TE}$ & Coeficientes de reflexión de Fresnel. \\
$f_c$, $B$, $\lambda_0$ & Frecuencia central (425\,MHz), ancho de banda (50\,MHz) y longitud de
                         onda en el aire. \\
$k_0$ & Número de onda en el aire, $2\pi/\lambda_0$. \\
$M$, $N_f$ & Número de pulsos y de frecuencias. \\
$Q_{img}$, $A_{PSF}$, $C_{voo}$, $J$ & Métricas del Exp.~F: reducción de ISLR por la fusión,
                                       crecimiento del volumen del lóbulo, costo de vuelo y costo
                                       ponderado total. \\
\bottomrule
\end{longtable}


% ══════════════════════════════════════════════════════════════════════════
\chapter{Introducción}
\label{cap:intro}

Este relatorio responde a las orientaciones del 03/10/2026. Tras el
relatorio del 28/09 quedó establecido que la potencia media recibida por el Rx
no varía sistemáticamente con su posición, ni frente al caso monoestático; la
optimización de potencia queda descartada. La pregunta pasa a ser
\emph{cuándo, dónde y por qué} un canal biestático adicional mejora la imagen
monoestática que el transceiver ya produce, y si operar cerca del
pseudo-Brewster suprime el retorno de la superficie lo suficiente para mejorar
la relación blanco--clutter. Las referencias de la forma §$n$ remiten a las
secciones de las Orientações del 03/10.

\section{Arquitectura}

El Tx es un \emph{transceiver}: transmite y recibe el canal monoestático
complejo en todos los pulsos. Un segundo UAV con el Rx remoto registra el
canal biestático \emph{del mismo pulso}. El producto relevante es la
comparación entre tres imágenes: monoestática, biestática y combinada
mono+bi.

\begin{table}[H]
\centering\small
\caption{Configuraciones de referencia (Orientações §1.1) tal como se simulan.}
\label{tab:config}
\begin{tabular}{@{}L{3.0cm}L{6.0cm}L{5.3cm}@{}}
\toprule
\textbf{Configuración} & \textbf{Definición en la simulación} & \textbf{Finalidad} \\
\midrule
Ref1 & Rx a la misma altura y radio, $\Delta\phi=180^{\circ}$ & baseline geométrico extremo (especular) \\
Ref2 & canal mono del transceiver ($\mathrm{Rx}=\mathrm{Tx}$) & referencia permanente \\
Bi intermedia & $\Delta\phi\in\{15^{\circ},\dots,165^{\circ}\}$ constante, $\kappa=1$ & diversidad angular \\
Multiestática & fusión coherente mono+bi (fus 2N y fusN con el mismo presupuesto) & información complementaria \\
Rx optimizado & $\Delta\phi(t)=\Delta\phi_0+A(2t/T-1)$, radio $\kappa\rho_{Tx}$, $z_{Rx}=z_{Tx}$ & Exp.~F \\
\bottomrule
\end{tabular}
\end{table}

\section{Preguntas de investigación (§4)}
\begin{enumerate}[itemsep=1pt]
  \item ¿Cómo altera la separación angular Tx--Rx el soporte en $k$ y la PSF?
  \item ¿Qué separación produce la mejor imagen biestática aislada?
  \item ¿Qué separación hace que el canal bi añada más información a la imagen mono?
  \item ¿Mejora la fusión coherente mono+bi la resolución horizontal, vertical, el volumen de la PSF, PSLR, ISLR, contraste o localización?
  \item ¿Se mantiene la mejora para objetos con RCS anisótropa y distintas orientaciones?
  \item ¿Reduce la operación cerca del pseudo-Brewster el retorno superficial de forma robusta (pérdidas, rugosidad, errores angulares)?
  \item ¿Cuál es el costo, en calidad de imagen, de imponer una restricción de pseudo-Brewster?
\end{enumerate}
Las respuestas, con las cifras que las sustentan, están en el
Capítulo~\ref{cap:respuestas}.

% ══════════════════════════════════════════════════════════════════════════
\chapter{Método}
\label{cap:metodo}

\section{Motor de simulación}

Se escribió un motor nuevo en Python/numba (\texttt{simulations/analise\_1610/bisar})
porque el estudio exige miles de imágenes 3D complejas con ecuación radar,
superficie y blancos PO; el núcleo de RCS ya estaba en Python
(\texttt{tools/rcs\_core}) y se portó la PO biestática de
\texttt{verifica\_ganancia\_brewster.m}. Cada pieza se validó contra el código
existente (Sección~\ref{sec:validacion}).

\paragraph{Geometría.} El Tx (transceiver) recorre la hélice óptima de
\texttt{run\_otimiza\_tx.m} (\texttt{radarTx\_espiral\_mono0.json}: 2 vueltas,
radio 137.4$\to$165.5\,m, altura 131.5$\to$93.6\,m, incidencia
$46.3^{\circ}$--$60.5^{\circ}$, media $53.4^{\circ}$), diezmada a $M=1600$ pulsos. El Rx remoto
vuela a la \emph{misma altura} que el Tx en cada instante,
\begin{equation}
  \vect{p}_{Rx}(t) = \big(\kappa(t)\rho_{Tx}(t)\cos[\phi_{Tx}(t)+\Delta\phi(t)],\;
  \kappa(t)\rho_{Tx}(t)\sin[\phi_{Tx}(t)+\Delta\phi(t)],\; z_{Tx}(t)\big),
\end{equation}
con $\kappa=1$ y $\Delta\phi$ constante en el Exp.~B (orientación §5.3) y
$\Delta\phi(t)$, $\kappa$ libres en el Exp.~F. Con $\kappa=1$ y $\Delta\phi$ constante, el
Rx recorre \emph{exactamente la misma hélice} que el Tx (mismos radios, mismas alturas,
misma velocidad), rotada rígidamente $\Delta\phi$ en acimut: es la hélice de
\texttt{funcao\_espiral} con \texttt{NorthOffset}$\,=\Delta\phi$, iniciada en el mismo instante
(Tx en acimut $0^{\circ}$, Rx en $\Delta\phi$). Ambos recorren sincronizados las dos vueltas, de
modo que la separación acimutal se mantiene constante en todos los pulsos. Se verificó
numéricamente que las dos construcciones coinciden ($<3\times10^{-13}$\,m). El blanco de referencia es un punto
isótropo de $\sigma=1\,\mathrm{m}^2$ en $(0,0,-2)$\,m.

\paragraph{Señal.} Portadora 425\,MHz, banda 50\,MHz muestreada en $N_f=64$
frecuencias (señal comprimida en rango), $P_t=100$\,kW, $G=6$\,dBi con haz
gaussiano $14^{\circ}\times50^{\circ}$ apuntado al centro, $T_p=1.3\,\mu$s, $T=300$\,K.
El eco del pulso $m$ en la frecuencia $f_k$ para el canal (Tx, Rx) es
\begin{equation}
  s_{mk} = \frac{\sqrt{P_t}\,\lambda_k}{(4\pi)^{3/2}}
  \frac{\sqrt{G_tG_r}}{R_TR_R}\,\tau_t\tau_r\,
  e^{-\alpha_k(R^{s}_T+R^{s}_R)}\,S_{mk}\,e^{-j2\pi f_k\tau_m},
  \qquad \tau_m=\frac{R^a_T+n'R^s_T+R^a_R+n'R^s_R}{c},
  \label{eq:eco}
\end{equation}
donde los tramos en aire ($R^a$) y suelo ($R^s$) salen del rayo de Fermat
(Newton sobre la ley de Snell), $\tau_{t,r}=\sqrt{1-|r_{TM/TE}|^2}$ son las
transmitancias de amplitud, $\alpha_k=k_0|n''|$ la atenuación y $S_{mk}$ la
amplitud compleja de dispersión ($S=1$ para el punto, PO para objetos). El
canal mono usa $\mathrm{Rx}=\mathrm{Tx}$ con \emph{el mismo pulso}.

\paragraph{Suelos.} $\varepsilon_r=\varepsilon'-j\varepsilon''$: S0 $=4$ (sin
pérdidas, validación), S1 $=3.99-j0.40$ (referencia de la orientación),
S2 $=6.5-j0.9$, S3 $=12-j2.5$. Pseudo-Brewster (mínimo de $|r_{TM}|^2$):
$63.43^{\circ}$, $63.45^{\circ}$, $68.66^{\circ}$ y $74.05^{\circ}$, con mínimos de
$-\infty$, $-34.5$, $-30.7$ y $-26.5$\,dB.

\paragraph{Superficie.} Dos componentes con la misma ecuación radar:
(i) \emph{especular coherente}: un punto especular por pulso (teoría de
imágenes, exacta por fase estacionaria para plano infinito), con
$r_{TM/TE}(\theta)$ y factor de Ament $\chi=\exp[-\tfrac12k_0^2s^2(\cos\theta_i+\cos\theta_s)^2]$;
(ii) \emph{difusa}: facetas de $0.15\times0.15$\,m en $\pm12$\,m con amplitud
$a_p\sqrt{\sigma^0_{pp}\,\Delta A}$, $a_p\sim\mathcal{CN}(0,1)$ \emph{congelada}
durante toda la apertura y compartida por los canales mono y bi, y
$\sigma^0$ del SPM de primer orden biestático (rugosidad moderada $s=2$\,cm,
$\ell=20$\,cm, $k_0s=0.18$):
\begin{equation}
  \sigma^0_{pp} = 8k_0^4s^2\cos^2\theta_i\cos^2\theta_s|\alpha_{pp}|^2
  W(k_0[\sin\theta_s\cos\phi_s-\sin\theta_i],\,k_0\sin\theta_s\sin\phi_s),
\end{equation}
\begin{equation}
  \alpha_{hh}=\frac{(\varepsilon-1)\cos\phi_s}{(\cos\theta_i+q_i)(\cos\theta_s+q_s)},\qquad
  \alpha_{vv}=\frac{(\varepsilon-1)\,[\varepsilon\sin\theta_i\sin\theta_s-\cos\phi_s\,q_iq_s]}
                  {(\varepsilon\cos\theta_i+q_i)(\varepsilon\cos\theta_s+q_s)},
  \label{eq:spm}
\end{equation}
$q=\sqrt{\varepsilon-\sin^2\theta}$, $\phi_s$ medido desde la dirección especular.
En retrodispersión $\alpha_{hh}=r_{TE}$; en dirección especular $\alpha_{vv}=0$
exactamente en Brewster. El error de atitud (roll $\gamma$) mezcla
polarizaciones: $r_{\mathrm{ef}}=\cos^2\gamma\,r_{TM}+\sin^2\gamma\,r_{TE}$.

\paragraph{Blancos.} PEC por PO biestática co-polar con integral exacta por
triángulo, polarización TM/TE general (vectores $\hat v,\hat h$ continuos en
la dirección, de modo que mono y bi convergen cuando $\Delta\phi\to0$),
evaluada en el suelo ($\lambda_s=\lambda_0/n'$) en las direcciones refractadas
y en cada frecuencia.

\paragraph{Imagen.} Back-projection con refracción exacta por vóxel y pulso:
\begin{equation}
  I(\vect r) = \frac{1}{\sum_m w_m N_f}\sum_m w_m\sum_k s_{mk}\,e^{\,j2\pi f_k\tau_m(\vect r)} ,
\end{equation}
(envolvente interpolada sobre el perfil sobremuestreado $32\times$, fase exacta).

\paragraph{Ponderación en rango y ventana de recepción.}\label{sec:ponderacion} Los Exp.~A, B, D y E
usan el filtro adaptado sin ponderar (espectro rectangular), igual que
\texttt{PROC\_snell\_test\_script.m}: es lo adecuado para la forma de la PSF.
En los experimentos limitados por clutter (C y F) se aplica además una
ponderación Hann en frecuencia (normalizada a media 1) y una ventana de
recepción de $[-30,+60]$\,m alrededor del camino Tx--centro--Rx. La razón es
un hallazgo del Exp.~C: en una órbita de altura constante, el reflejo
especular de nadir del canal mono tiene retardo constante, y sus lóbulos
lejanos en rango ($\approx-38$\,dB a 25 celdas con espectro rectangular) se
acumulan coherentemente sobre el eje de la órbita, justo donde está el blanco.
Con Hann la fuga cae de $-78$ a $-129$\,dBm (pico del blanco $-69.08$\,dBm,
pérdida de 0.08\,dB), a cambio de un lóbulo vertical $\approx1.6$ veces más
ancho. El Exp.~C se repitió con espectro rectangular como análisis de
sensibilidad (\texttt{io/analise\_1610/expC\_rect}) y sus conclusiones no cambian.

\section{Normalización de la fusión}
\label{sec:normalizacion}

Para que el aumento del número de muestras no se confunda con ganancia
geométrica (§5.1, §11), cada imagen es la \emph{media} coherente de sus muestras y
se definen dos fusiones:
\begin{itemize}[itemsep=2pt]
  \item \textbf{fus (2N)}: $I_{fus}=\tfrac12(I_{mono}+I_{bi})$, ambos canales completos
  (lo que el sistema adquiere realmente).
  \item \textbf{fusN (N, mismo presupuesto)}: canal mono en los pulsos impares y
  bi en los pares, con el mismo número total de muestras que la imagen mono.
\end{itemize}
Las métricas de forma (resoluciones $\delta_{x,y,z}$, volumen, anisotropía, PSLR, ISLR, contraste) son
invariantes a la escala; la SNR se calcula propagando el ruido térmico,
$\mathrm{var}=kT/(T_pM_{\mathrm{ef}})$, con $M_{\mathrm{ef}}=M$ (mono, bi, fusN)
y $2M$ (fus). Así el $+3$\,dB por duplicar muestras aparece \emph{sólo} en la
SNR de fus y se declara como tal. Todas las imágenes se muestran con la misma
referencia (pico monoestático $=0$\,dB) y el mismo rango dinámico (30\,dB).

\section{Métricas}

\begin{table}[H]
\centering\small
\caption{Métricas implementadas (\texttt{bisar/metrics.py}).}
\label{tab:metricas}
\begin{tabular}{@{}L{3.1cm}L{11.2cm}@{}}
\toprule
$\delta_{x,y,z}$ & ancho a $-3$\,dB en cortes por el pico (interpolación lineal en dB);
               censurado si no cruza dentro de la ventana. Equivale al FWHM en potencia
               (columnas \texttt{fwhm\_*} de los CSV); formas cerradas en el Apéndice~\ref{ap:resolucion} \\
$V_{-3\,\mathrm{dB}}$ & volumen de la componente conexa $|I|^2\ge\tfrac12$pico que contiene el pico \\
Anisotropía & $\sqrt{\lambda_{max}/\lambda_{min}}$ del tensor de segundo momento del lóbulo $-3$\,dB (3D y $xy$) \\
PSLR, ISLR & lóbulo principal = elipsoide con semiejes al primer nulo en $x,y,z$; ventana fija
             $2\times2\times5$\,m; ISLR = energía fuera / dentro \\
Contraste & pico / potencia media fuera del lóbulo principal \\
Localización & distancia del pico (refinamiento parabólico) a la posición verdadera \\
Amplitud & $20\log_{10}(|I_{pico}|/\overline{|a|})$ (pérdida de enfoque) \\
Separabilidad & fracción de blancos con máximo local y caída $\ge3$\,dB hacia cada vecino \\
Dominio $k$ & $\vect K=\tfrac{2\pi fn'}{c}(\hat u_T+\hat u_R)$; área ocupada en $(k_x,k_y)$ (celdas de $0.25$\,rad/m),
              dispersión $\sigma_{k}$ (Fisher) \\
Condicionamiento & matriz $A_{iv}=e^{-j\vect K_i\cdot\vect r_v}$ sobre $7\times7\times5$ vóxeles; rango efectivo
                   (Roy--Vetterli), coherencia mutua; Gram normalizado por canal \\
SCR & $|I_{blanco}(\vect r_0)|^2/\mathrm{E}|I_{sup}|^2$ en la vecindad del blanco (imagen) y
      $\bar P_{blanco}/\bar P_{sup}$ (enlace) \\
SNR & $|I_{pico}|^2T_pM_{\mathrm{ef}}/kT$ \\
Vuelo & separación mínima Tx--Rx, $v_{max}^{Rx}/v_{Tx}$, longitud relativa, aceleración a $v_{Tx}=10$\,m/s \\
\bottomrule
\end{tabular}
\end{table}

\section{Validación del motor}
\label{sec:validacion}

\begin{table}[H]
\centering\small
\caption{Validaciones (\texttt{analise\_1610/tests}).}
\label{tab:validacion}
\begin{tabular}{@{}L{5.6cm}L{4.4cm}L{4.4cm}@{}}
\toprule
\textbf{Prueba} & \textbf{Referencia} & \textbf{Motor nuevo} \\
\midrule
RCS PO en nadir (Tabla 2 del 28/09): cubo, caja, cilindro, esfera, cono
 & 20.04, 14.02, 8.94, $-1.54$, $-15.86$\,dBsm & idénticos (0.00\,dB) \\
Placa especular $4\pi A^2\cos^2\theta/\lambda^2$, $\theta=15,30,45^{\circ}$
 & 19.74, 18.79, 17.03\,dBsm & 19.73, 18.88, 17.01\,dBsm \\
BP vs.\ núcleo de PROC (MATLAB), 10\,GHz, corte fino $x$ & 6.66\,mm & 6.71\,mm \\
BP vs.\ PROC, eje $z$ (mono0 / offset0) & 0.200 / 0.206\,m & 0.202 / 0.207\,m \\
PSF horizontal mono a 425\,MHz vs.\ anillo de Bessel & 0.157\,m & 0.1575\,m \\
PSF bi $\Delta\phi=90^{\circ}$ vs.\ $\delta_{xy}/\cos(\Delta\phi/2)$ & 0.222\,m & 0.221\,m \\
Suma de facetas de Kirchhoff (plano liso, Ref1) vs.\ teoría de imágenes & $-0.102$\,dBm & $-0.095$\,dBm \\
Potencia difusa media (40 realizaciones) vs.\ $\int\sigma^0dA$ & 0\,dB & $-0.24$ a $+0.45$\,dB \\
SPM: $\alpha_{hh}$ en retrodispersión vs.\ $r_{TE}$; $\alpha_{vv}$ especular en Brewster
 & $-0.5104$; 0 & $-0.5104$; 0 \\
\bottomrule
\end{tabular}
\end{table}

% ══════════════════════════════════════════════════════════════════════════
\chapter{Experimento A: referencias de imagen}
\label{cap:resA}

\paragraph{Hipótesis.} Ref2 (monoestática) es la referencia permanente; Ref1
(diametral, $\Delta\phi=180^{\circ}$, misma altura y radio) es un \emph{baseline}
geométrico extremo que debe reconstruir mal; su fusión con Ref2 no debería
mejorar la imagen.

\paragraph{Método.} Punto isótropo de $1\,\mathrm{m}^2$ en $(0,0,-2)$\,m, suelo
S1, mismos pulsos para los dos canales; rejilla $2\times2\times5$\,m con
$\Delta=2.5$\,cm/10\,cm; cortes de $\pm25$\,m para medir anchos censurados.
Datos complejos en \texttt{expA/datos\_complejos\_expA.npz}.

\begin{figure}[H]
\centering
\includegraphics[width=\textwidth]{A_geometria.pdf}
\caption{Geometría completa: Tx transceiver (canal
mono = Ref2), Rx remoto en Ref1 ($\Delta\phi=180^{\circ}$) y en una posición
intermedia ($\Delta\phi=90^{\circ}$); a la derecha, el perfil radial común de las
hélices y las líneas de incidencia (mín., pseudo-Brewster de S1, máx.).}
\label{fig:geom}
\end{figure}

\begin{table}[H]
\centering\small
\caption{Exp.~A: métricas de la PSF (\texttt{expA\_metricas.csv}). $\delta_x^{\pm25}$: ancho en el corte largo.}
\label{tab:expA}
\begin{tabular}{@{}lrrrrrrrrr@{}}
\toprule
Canal & $\delta_x$ & $\delta_x^{\pm25}$ & $\delta_z$ & $V_{-3}$ & Aniso & PSLR & ISLR & Contr. & SNR \\
      & [m] & [m] & [m] & [l] & & [dB] & [dB] & [dB] & [dB] \\
\midrule
Ref2 mono              & 0.158 & 0.159 & 1.39 & 19.2 & 8.9 & $-7.9$ & 6.4 & 21.3 & 78.0 \\
Ref1 bi (180°)         & $>2$ & \textbf{17.2} & 1.39 & $>5331$ & --- & $0.0$ & --- & 8.5 & 78.0 \\
Ref1+Ref2 (2N)         & 0.235 & --- & 1.39 & 41.1 & 5.9 & $-3.8$ & 11.9 & 11.4 & 81.0 \\
Ref1+Ref2 (N)          & 0.235 & --- & 1.39 & 41.1 & 5.9 & $-3.8$ & 11.9 & 11.4 & 78.0 \\
\midrule
Bi intermedio 90°      & 0.221 & 0.223 & 1.39 & 38.8 & 6.2 & $-7.9$ & 5.1 & 19.6 & 78.0 \\
Mono+bi 90° (N)        & 0.182 & --- & 1.39 & 25.9 & 7.4 & $-9.9$ & 1.0 & 25.3 & 78.0 \\
\bottomrule
\end{tabular}
\end{table}

\begin{figure}[H]
\centering
\includegraphics[width=\textwidth]{A_psf_cortes.pdf}
\caption{Cortes de la PSF en $x$, $y$ y $z$ para Ref2, Ref1 y su fusión (misma
escala: pico mono $=0$\,dB). Ref1 es plana en el plano horizontal: no enfoca.}
\label{fig:A_cortes}
\end{figure}

\begin{figure}[H]
\centering
\includegraphics[width=\textwidth]{A_imagenes.pdf}
\caption{Imágenes de referencia con la misma escala dinámica (30\,dB):
plano $z=-2$\,m y corte $y=0$.}
\label{fig:A_img}
\end{figure}

\begin{figure}[H]
\centering
\includegraphics[width=0.85\textwidth]{A_kspace_ref.pdf}
\caption{Dónde se pierde el soporte en Ref1: el canal mono cubre un anillo
horizontal de radio $2k_0\sin\theta_i\approx14$\,rad/m; en Ref1 las
direcciones de ida y vuelta en el suelo son simétricas respecto de la vertical,
$\vect K_h=k_sn'(\hat u_{T,h}+\hat u_{R,h})=0$ para todo pulso, y todo el
soporte cae sobre el eje $k_z$.}
\label{fig:A_k}
\end{figure}

\paragraph{Resultados e interpretación.}
\begin{itemize}[itemsep=2pt]
  \item La potencia media recibida es idéntica en Ref1 y Ref2
  ($-68.87$\,dBm), confirmando la conclusión del 28/09 en esta geometría.
  \item \textbf{Ref1 reconstruye mal, como se esperaba.} Su soporte en $k$ es
  un segmento sobre el eje $k_z$ (Fig.~\ref{fig:A_k}); conserva exactamente la
  resolución vertical (1.39\,m) pero pierde toda la horizontal: el único
  enfoque lateral residual proviene de la curvatura del frente de onda
  ($\delta_x=17$\,m con $R\approx190$\,m).
  \item \textbf{La fusión Ref1+Ref2 es peor que Ref2}: el canal Ref1 añade un
  pedestal plano a la PSF, sube los lóbulos (PSLR $-7.9\to-3.8$\,dB, ISLR
  $+6.4\to+11.9$\,dB) y ensancha el lóbulo principal en un 49\,\%. Es un
  resultado negativo claro: la geometría diametral no debe usarse.
  \item La fusión con el mismo presupuesto (N) da exactamente la misma imagen
  que la de 2N (correlación $>0.99999$): todo cambio de forma es geométrico.
  La única diferencia es la SNR ($+3.01$\,dB), atribuible solo al número de
  muestras.
\end{itemize}

\conclexp{A}{Ref1 conserva la resolución vertical pero pierde toda la horizontal
($\delta_x=17$\,m) y su fusión empeora la mono (PSLR $-7.9\to-3.8$\,dB), de modo que sólo
sirve como referencia geométrica extrema y no como configuración de imagen.}

% ══════════════════════════════════════════════════════════════════════════
\chapter{Experimento B: barrido de la separación Tx--Rx}
\label{cap:resB}

\paragraph{Hipótesis.} El canal bi a la misma altura y radio contrae el anillo
horizontal de soporte por un factor $|\cos(\Delta\phi/2)|$; un anillo interior
complementario podría reducir los lóbulos del anillo monoestático (función
$J_0$, PSLR $\approx-7.9$\,dB) a costa de ensanchar el lóbulo principal.

\paragraph{Método.} $\Delta\phi=0^{\circ}{:}15^{\circ}{:}180^{\circ}$ constante, $z_{Rx}=z_{Tx}$,
$\kappa=1$; cuatro imágenes por caso (mono, bi, fus, fusN), blanco en el eje
y blanco descentrado en $(6,4,-2)$\,m; soporte en $k$ y condicionamiento.
Barrido directo, sin optimizador (§5.3). CSV: \texttt{expB\_metricas.csv},
\texttt{expB\_kspace\_condicionamiento.csv}.

\begin{figure}[H]
\centering
\includegraphics[width=\textwidth]{B_geometria.pdf}
\caption{Geometría del Exp.~B. Arriba, planta del Tx (azul) y del Rx (naranja) para cada
$\Delta\phi$: el Rx recorre \emph{la misma hélice} que el Tx, con los mismos radios y alturas,
rotada $\Delta\phi$ en acimut y sincronizada (marcas = mismos pulsos; línea punteada = línea
de base). Abajo, $\rho(t)$ y $z(t)$ de ambos, idénticos en todos los casos. Lo único que
cambia entre columnas es el desfase acimutal. Las dos espiras se ven separadas en planta
porque, en un mismo acimut, Tx y Rx pasan en instantes distintos y, por tanto, con radios
distintos de la hélice.}
\label{fig:B_geo}
\end{figure}

\paragraph{Lectura de las figuras siguientes.} Los anillos de las Figs.~\ref{fig:B_k} y
\ref{fig:B_img} \emph{no son trayectorias}. En la Fig.~\ref{fig:B_k}, cada punto es el
vector de onda biestático $\vect{k}=k(\hat{\vect u}_{Tx}+\hat{\vect u}_{Rx})$ de un pulso, en
rad/m. Como los dos UAV vuelan a la misma altura y radio, su componente horizontal vale
$2k\sin\theta\,|\cos(\Delta\phi/2)|$: el anillo bi se contrae con $\Delta\phi$ aunque la
órbita del Rx sea la misma. En la Fig.~\ref{fig:B_img}, los anillos son los lóbulos de la
PSF (tipo $J_0$), en metros alrededor del blanco.

\begin{figure}[H]
\centering
\includegraphics[width=\textwidth]{B_kspace.pdf}
\caption{Soporte en el dominio de $k$: Ref2, tres separaciones
intermedias y Ref1. Arriba, plano $(k_x,k_y)$; abajo, $(|k_h|,k_z)$. El canal
bi es un anillo de radio $|\cos(\Delta\phi/2)|$ veces el monoestático, con la
misma extensión en $k_z$. Son números de onda (rad/m), no trayectorias: la órbita del Rx
es la misma en todas las columnas (Fig.~\ref{fig:B_geo}).}
\label{fig:B_k}
\end{figure}

\begin{figure}[H]
\centering
\includegraphics[width=\textwidth]{B_resoluciones.pdf}
\caption{Resoluciones $x$, $y$, $z$ vs.\ $\Delta\phi$. Curva
discontinua: $\delta_{xy}^{mono}/|\cos(\Delta\phi/2)|$ (forma cerrada, ec.~\ref{eq:A_dxy_eje}).}
\label{fig:B_res}
\end{figure}

\begin{figure}[H]
\centering
\includegraphics[width=0.9\textwidth]{B_volumen_anisotropia.pdf}
\caption{Volumen del lóbulo $-3$\,dB y anisotropía 3D de la PSF vs.\ $\Delta\phi$.}
\label{fig:B_vol}
\end{figure}

\begin{figure}[H]
\centering
\includegraphics[width=\textwidth]{B_pslr_islr.pdf}
\caption{PSLR, ISLR y contraste vs.\ $\Delta\phi$. La línea azul es la
imagen monoestática; verde, la fusión; punteado, la fusión con el mismo
presupuesto de muestras (coincide).}
\label{fig:B_lob}
\end{figure}

\begin{figure}[H]
\centering
\includegraphics[width=\textwidth]{B_informacion.pdf}
\caption{Información: área ocupada en $(k_x,k_y)$, rango efectivo y coherencia
mutua de la matriz de adquisición linealizada ($7\times7\times5$ vóxeles,
Gram normalizado por canal).}
\label{fig:B_info}
\end{figure}

\begin{figure}[H]
\centering
\includegraphics[width=\textwidth]{B_imagenes.pdf}
\caption{Imágenes lado a lado con la misma escala (pico mono $=0$\,dB, 30\,dB):
mono, bi y fusión para $\Delta\phi=45^{\circ},90^{\circ},135^{\circ},180^{\circ}$. Los anillos son
los lóbulos secundarios de la PSF en una ventana de $2\times2$\,m alrededor del blanco.}
\label{fig:B_img}
\end{figure}

\begin{table}[H]
\centering\small
\caption{Exp.~B, blanco en el eje: canal bi aislado y fusión (N) frente a mono.
Mono: $\delta_x=0.157$\,m, $V=19.2$\,l, PSLR $-7.88$\,dB, ISLR $6.37$\,dB,
contraste 21.3\,dB, rango efectivo 28.4.}
\label{tab:expB}
\begin{tabular}{@{}r|rrr|rrrrrr@{}}
\toprule
 & \multicolumn{3}{c|}{\textbf{Bi aislado}} & \multicolumn{6}{c}{\textbf{Fusión mono+bi (N)}} \\
$\Delta\phi$ & $\delta_x$ & $V$ & ISLR & $\delta_x$ & $V$ & PSLR & ISLR & Contr. & rango ef. \\
\midrule
 15 & 0.158 & 19.7 & 6.36 & 0.158 & 19.6 & $-7.90$ & 6.35 & 21.3 & 28.3 \\
 30 & 0.162 & 20.8 & 6.31 & 0.160 & 20.2 & $-7.92$ & 6.21 & 21.4 & 28.1 \\
 45 & 0.169 & 22.6 & 6.19 & 0.163 & 21.0 & $-8.00$ & 5.71 & 21.7 & 28.2 \\
 60 & 0.180 & 25.6 & 5.87 & 0.167 & 21.9 & $-8.28$ & 4.54 & 22.6 & 28.9 \\
 75 & 0.199 & 29.4 & 5.58 & 0.174 & 23.4 & $-8.79$ & 2.53 & 24.2 & 30.4 \\
 90 & 0.221 & 38.8 & 5.10 & 0.182 & 25.9 & $-9.91$ & \textbf{1.04} & \textbf{25.3} & 32.4 \\
105 & 0.258 & 52.0 & 4.37 & 0.192 & 28.0 & $-12.19$ & 1.60 & 24.3 & \textbf{33.8} \\
120 & 0.315 & 75.9 & 3.45 & 0.204 & 30.3 & $\mathbf{-13.09}$ & 2.50 & 22.8 & 33.7 \\
135 & 0.412 & 131 & 1.62 & 0.214 & 34.8 & $-10.24$ & 2.32 & 22.2 & 31.5 \\
150 & 0.612 & 285 & $-1.13$ & 0.224 & 37.6 & $-9.48$ & 3.47 & 20.4 & 26.7 \\
165 & 1.216 & 1118 & $-9.6$ & 0.232 & 40.8 & $-5.25$ & 7.30 & 16.2 & 20.0 \\
180 & 17.2 & $>5331$ & --- & 0.235 & 41.1 & $-3.81$ & 11.95 & 11.4 & 16.0 \\
\bottomrule
\end{tabular}
\end{table}

\paragraph{Resultados e interpretación.}
\begin{enumerate}[itemsep=3pt]
  \item \textbf{Ley geométrica verificada.} La resolución horizontal del canal bi
  sigue $\delta_{xy}^{mono}/|\cos(\Delta\phi/2)|$ con error $<1$\,\% hasta
  150°; la vertical no cambia (1.386\,m en todos los casos) porque, a la misma
  altura, la extensión $k_z$ es idéntica ($\sigma_{kz}=1.174$\,rad/m constante).
  \item \textbf{Mejor imagen biestática aislada (pregunta 2):} la de menor separación
  viable. Toda separación degrada el lóbulo principal del bi aislado; no
  existe un óptimo intermedio para el bi solo. Con una separación mínima de
  seguridad de 15\,m entre UAVs ($\Delta\phi\gtrsim6^{\circ}$), el bi aislado más
  nítido es prácticamente un segundo monoestático y no aporta información.
  \item \textbf{Mayor complementariedad (pregunta 3): $\Delta\phi\approx90^{\circ}$--$120^{\circ}$.} El anillo
  interior rellena el disco de $k$: el área ocupada en $(k_x,k_y)$ crece de
  294 a 475\,rad$^2$/m$^2$ (75°) y el rango efectivo de 28.4 a 33.8 (105°).
  En la imagen fusionada, el ISLR cae 5.3\,dB (90°), el PSLR mejora 5.2\,dB
  (120°) y el contraste sube 4.0\,dB (90°).
  \item \textbf{Costo:} el lóbulo principal horizontal se ensancha
  ($\delta_x$ $+16$\,\% a 90°, $+30$\,\% a 120°; volumen $+35$\,\% y $+58$\,\%).
  La fusión no mejora la resolución $-3$\,dB: mejora los lóbulos, el contraste y el
  condicionamiento. La anisotropía 3D baja (8.9$\to$7.4 a 90°) sólo porque el lóbulo
  horizontal se ensancha, no porque mejore $\delta_z$.
  \item \textbf{Resultados idénticos para el blanco descentrado} ($6,4,-2$\,m;
  diferencias $<0.3$\,dB en PSLR): con la hélice de dos vueltas completas la
  PSF es casi invariante en la escena.
  \item \textbf{Normalización:} fus y fusN coinciden con correlación $1.00000$ en
  los 26 casos. Ninguna mejora de forma se debe al doble de muestras.
\end{enumerate}

\conclexp{B}{El canal bi a la misma altura no mejora ninguna resolución (la horizontal empeora
como $1/|\cos(\Delta\phi/2)|$ y la vertical no cambia), pero con $\Delta\phi\approx90^{\circ}$--$120^{\circ}$
rellena el dominio $k$ y la fusión baja el ISLR 5.3\,dB, mejora el PSLR hasta 5.2\,dB y sube el
contraste 4\,dB, a cambio de un lóbulo principal un 16--30\,\% más ancho.}

% ══════════════════════════════════════════════════════════════════════════
\chapter{Experimento C: incidencia con y sin pseudo-Brewster}
\label{cap:resC}

\paragraph{Hipótesis (§2.1).} Cerca del pseudo-Brewster se \emph{reduce} el
retorno de la interfaz en TM (no aumenta el eco del blanco), lo que mejoraría
la relación blanco--superficie si esa reducción es robusta a pérdidas,
rugosidad, errores angulares y de atitud.

\paragraph{Método.} Se comparan dos familias de trayectoria del Tx (F1, hélice; F2,
incidencia constante), descritas abajo, en varias incidencias, suelos, rugosidades,
polarizaciones y errores de apunte (Tabla~\ref{tab:C_var}). Canales: mono, bi a 90°, 135°
y 180° (Ref1), y fusión. El SCR se mide en el vóxel del blanco (punto de
$1\,\mathrm{m}^2$ a 2\,m de profundidad), con la superficie congelada y compartida por los
dos canales (3 realizaciones), ponderación Hann y ventana de recepción
(Sección~\ref{sec:ponderacion}). CSV:
\texttt{expC\_enlace\_reflectancia.csv}, \texttt{expC\_enlace\_error\_roll.csv},
\texttt{expC\_scr\_imagen.csv}, \texttt{expC\_psf\_vs\_incidencia.csv}; versión
con espectro rectangular en \texttt{expC\_rect/}.

\paragraph{Geometría de F1 y F2.} Las dos familias se construyen punto a punto a partir de
la hélice base (\texttt{geometry.familia\_C}). Cada pulso conserva su acimut y su distancia al
centro de la escena, que en la hélice óptima es casi constante ($\approx189$\,m, porque su línea
de base es perpendicular a la línea de vista). Sólo cambia la incidencia $\theta$, medida en el
centro de la superficie $(0,0,0)$:
\begin{itemize}[itemsep=1pt]
  \item \textbf{F1:} cada punto se gira en su plano vertical $\theta_c-\langle\theta_0\rangle$. La
  trayectoria sigue siendo una hélice y conserva su barrido propio de incidencia
  ($\theta_c\pm7^{\circ}$).
  \item \textbf{F2:} cada punto se coloca exactamente en $\theta_c$. Como la distancia es casi
  constante, F2 es una órbita de dos vueltas a altura y radio constantes; en
  $\theta_{pB}=63.45^{\circ}$, 84.6\,m de altura y 169\,m de radio.
\end{itemize}
Un \emph{cono de incidencia constante} $\theta$ es el conjunto de puntos que, vistos desde
el centro de la superficie, forman el mismo ángulo $\theta$ con la vertical: un cono con vértice
en $(0,0,0)$ y semiángulo $\theta$ (Fig.~\ref{fig:C_3d}). La hélice base y F1 lo atraviesan
(su incidencia cambia en cada pulso); F2 vive sobre él.

\begin{figure}[H]
\centering
\includegraphics[width=\textwidth]{C_trayectorias_3d.pdf}
\caption{Trayectorias completas del Tx en 3D (\texttt{exp\_C\_fig\_trayectorias3d.py}). (a)
Hélice base, entre los conos de $46.3^{\circ}$ y $60.5^{\circ}$. (b) F1 en $\theta_{pB}$: la
misma hélice girada hacia abajo, que cruza el cono de $\theta_{pB}$ (rojo) a mitad de la
apertura. (c) F2 en $\theta_{pB}$: todos los pulsos sobre el cono de $63.45^{\circ}$, es decir,
un círculo de 169\,m de radio a 85\,m de altura recorrido dos veces.}
\label{fig:C_3d}
\end{figure}

El Rx sigue a su Tx a la misma altura y el mismo radio, desfasado $\Delta\phi$, como en el
Exp.~B. En F2 ve, por tanto, la misma incidencia $\theta_c$.

La reflexión especular no ocurre en $\theta_c$, sino en el punto especular del par Tx--Rx, que
está a mitad de camino entre sus nadires (Fig.~\ref{fig:C_geo}\,c). Su incidencia es
\begin{equation}
  \theta_{esp}=\arctan\!\big(\tan\theta_c\,\sin(\Delta\phi/2)\big),
  \label{eq:C_esp}
\end{equation}
y está a $\rho\,|\cos(\Delta\phi/2)|$ del centro. En el canal mono la especular es el reflejo de
nadir ($\theta_{esp}=0$); a $90^{\circ}$, con $\theta_c=\theta_{pB}$, $\theta_{esp}=54.8^{\circ}$; a
$135^{\circ}$, $61.6^{\circ}$. Sólo en Ref1 ($180^{\circ}$) coincide con $\theta_c$
(Fig.~\ref{fig:C_geo}\,d). \textbf{Por eso, elegir $\theta_c=\theta_{pB}$ sólo pone la
reflexión especular en Brewster en Ref1.} En las demás configuraciones, lo que se suprime en
$\theta_c$ es el retorno difuso. El modelo calcula la especular en el punto especular real de
cada pulso (\texttt{scene.specular\_echo}, plano infinito), no en $\theta_c$.

\begin{figure}[H]
\centering
\includegraphics[width=\textwidth]{C_geometria.pdf}
\caption{Geometría del Exp.~C (\texttt{exp\_C\_fig\_geometria.py}). (a) Plano vertical
$(\rho,z)$: hélice base, F1 en $\theta_{pB}$ (segmento de $\pm7^{\circ}$) y F2 (un punto por
$\theta_c$, con los errores $\theta_{pB}+\delta$ y las incidencias fuera de Brewster). (b)
Incidencia del Tx en cada pulso. (c) Planta de F2 en $\theta_{pB}$, con el Rx a 90°, 135° y
180° y el punto especular de cada par. (d) Incidencia en el punto especular frente a
$\Delta\phi$: ecuación (\ref{eq:C_esp}) y valores de la simulación.}
\label{fig:C_geo}
\end{figure}

\begin{table}[H]
\centering\small
\caption{Variables del Exp.~C y cómo entran en la simulación (\texttt{params/expC.json},
\texttt{exp\_C\_brewster.py}).}
\label{tab:C_var}
\begin{tabular}{@{}L{2.5cm}L{3.3cm}L{8.6cm}@{}}
\toprule
\textbf{Variable} & \textbf{Valores} & \textbf{Qué cambia en la simulación} \\
\midrule
Familia & F1, F2 & Forma de la trayectoria del Tx (y del Rx): hélice con barrido de $\pm7^{\circ}$
o cono de incidencia constante \\
$\theta_c$ & 35°, 45°, 53.4° (base), 75°; $\theta_{pB}+\delta$ & Incidencia central del Tx en el
centro de la superficie; fija la altura y el radio de la órbita \\
$\theta_{pB}$ & S0 63.43°, S1 63.45°, S2 68.66°, S3 74.05° & Mínimo de $|r_{TM}(\theta)|^2$ del suelo;
sólo sirve para elegir $\theta_c$ \\
$\delta$ (error angular) & 0, $\pm0.5$, $\pm1$, $\pm2$, $\pm5$° & Desplaza toda la trayectoria:
$\theta_c=\theta_{pB}+\delta$ (error sistemático de apunte o de altura) \\
$\Delta\phi$ & 0 (mono), 90°, 135°, 180° & Rx en la misma órbita desfasado $\Delta\phi$; fija
$\theta_{esp}$ (ec.~\ref{eq:C_esp}) y la posición del punto especular \\
Polarización & TM, TE & Coeficiente de Fresnel en la especular, núcleo SPM ($\approx$VV o HH) en la
difusa y transmitancia del blanco; Tx y Rx co-polares \\
Roll $\gamma$ & 0, 1, 2, 5° & Antena girada alrededor de la línea de vista: pesos
$\cos^2\gamma$ (TM) y $\sin^2\gamma$ (TE); sólo F2 en $\theta_{pB}$, S1, TM \\
Suelo & S0--S3 (S0 4, S1 3.99$-j$0.40, S2 6.5$-j$0.9, S3 12$-j$2.5) & Fresnel, $\sigma^0$ SPM,
transmitancia y $\theta_{pB}$. S1 completo; S0, S2 y S3 en $\theta_{pB}$, $\theta_{pB}\pm2^{\circ}$ y 53.4° \\
Rugosidad $s$ & 0, 1, 2, 4\,cm ($\ell=20$\,cm) & Especular $\times\,e^{-2k_0^2s^2\cos^2\theta_{esp}}$
(componente coherente); difusa SPM con amplitud $\propto s$ sobre un parche de $\pm12$\,m
(facetas de 15\,cm, 3 realizaciones) \\
SCR & enlace e imagen & Enlace: potencias $P_{blanco}/(P_{esp}+P_{dif})$. Imagen:
$|I_{blanco}|^2/\langle|I_{sup}|^2\rangle$ en el vóxel del blanco tras la back-projection
(ventana Hann y ventana de recepción) \\
\bottomrule
\end{tabular}
\end{table}

\section{Enlace: reflectancia y su robustez}

\begin{figure}[H]
\centering
\includegraphics[width=\textwidth]{C_reflectancia.pdf}
\caption{Reflectancia TM y TE de los suelos (izq.) y reflectancia TM
efectiva frente al error de incidencia y de atitud para S0 (sin pérdidas) y
S1 ($\varepsilon_r=3.99-j0.40$).}
\label{fig:C_refl}
\end{figure}

\begin{table}[H]
\centering\small
\caption{Reflectancia TM efectiva [dB] en $\theta_{pB}+\delta$ (\texttt{expC\_enlace\_error\_roll.csv}).
Roll 0°: media de $\pm\delta$ (simétrica a $<1$\,dB); roll 5°: pares $+\delta$\,/\,$-\delta$.
Última columna: incidencia base 53.4°. Menor es mejor (más supresión). Color según la supresión
respecto de la incidencia base (en el peor de los dos signos de $\delta$):
\colorbox{mejor}{$\ge10$\,dB}, \colorbox{regular}{5--10\,dB}, \colorbox{peor}{$<5$\,dB}.}
\label{tab:C_err}
\resizebox{\textwidth}{!}{\begin{tabular}{@{}lrrrrrr|r@{}}
\toprule
Suelo / roll & $\delta=0$ & $\pm0.5^{\circ}$ & $\pm1^{\circ}$ & $\pm2^{\circ}$ & $\pm5^{\circ}$ & & fuera (53.4°) \\
\midrule
S0, roll 0° & \cellcolor{mejor}$-189.7$ & \cellcolor{mejor}$-41.7$ & \cellcolor{mejor}$-35.7$ & \cellcolor{mejor}$-29.7$ & \cellcolor{peor}$-21.7$ & & $-17.5$ \\
S0, roll 5° & \cellcolor{mejor}$-46.8$ & \cellcolor{mejor}$-37.8$/$-49.1$ & \cellcolor{mejor}$-33.4$/$-38.9$ & \cellcolor{mejor}$-28.2$/$-31.5$ & \cellcolor{peor}$-20.3$/$-23.3$ & & --- \\
S1, roll 0° & \cellcolor{mejor}$-34.5$ & \cellcolor{mejor}$-33.8$ & \cellcolor{mejor}$-32.1$ & \cellcolor{mejor}$-28.5$ & \cellcolor{peor}$-21.5$ & & $-17.5$ \\
S1, roll 5° & \cellcolor{mejor}$-34.4$ & \cellcolor{mejor}$-32.9$/$-34.5$ & \cellcolor{mejor}$-31.0$/$-33.3$ & \cellcolor{regular}$-27.3$/$-29.8$ & \cellcolor{peor}$-20.1$/$-23.0$ & & --- \\
S2 ($\theta_{pB}=68.66^{\circ}$) & \cellcolor{mejor}$-30.7$ & & & & & & $-12.6$ \\
S3 ($\theta_{pB}=74.05^{\circ}$) & \cellcolor{mejor}$-26.5$ & & & & & & $-8.7$ \\
\bottomrule
\end{tabular}}
\end{table}

\begin{itemize}[itemsep=2pt]
  \item En el suelo sin pérdidas el mínimo es un cero exacto, pero \textbf{desaparece
  con 0.5°} de error ($-190\to-42$\,dB): es el ``mínimo extremadamente estrecho'' que
  §2.1 advierte que no basta para justificar una trayectoria.
  \item Con pérdidas el mínimo es \textbf{poco profundo pero ancho}: en S1 sólo
  $-34.5$\,dB, y se pierden 0.7, 2.4, 6.0 y 13.0\,dB con errores de 0.5°, 1°, 2° y 5°.
  Respecto de la incidencia base (53.4°), la supresión disponible es de 17.0\,dB.
  El mínimo empeora con la humedad ($-30.7$ y $-26.5$\,dB en S2 y S3) y se desplaza
  a 68.7° y 74.0°.
  \item El roll de 1°--5° es irrelevante en S1 (variación $\le1.5$\,dB); en S0 sólo
  desplaza el cero $\approx0.3^{\circ}$ (modelo co-polar, sin término cruzado).
\end{itemize}

\paragraph{Brewster biestático del clutter difuso.} La ec.~\eqref{eq:spm} muestra
que, para el retorno difuso, el pseudo-Brewster no es un punto sino una \emph{curva} en el
espacio (incidencia, separación): para $\theta<\theta_{pB}$ existe una separación
$\Delta\phi^\ast(\theta)<180^{\circ}$ donde el retorno difuso VV de primer orden se anula.
La curva termina en el Brewster clásico ($\Delta\phi=180^{\circ}$, $\theta=\theta_{pB}$) y
nunca pasa por la mono ($\Delta\phi=0$). En HH el nulo está en $\Delta\phi=90^{\circ}$ para
cualquier incidencia y es puramente geométrico. Comparar TM con TE en la misma geometría
separa ambos efectos (§5.4).

\begin{figure}[H]
\centering
\includegraphics[width=\textwidth]{C_sigma0_mapa_suelos.pdf}
\caption{$\sigma^0$ biestático SPM con el Rx a la misma elevación que el Tx, para los
cuatro suelos ($s=2$\,cm). Arriba VV (TM), abajo HH (TE). Curva roja: mínimo de VV en cada
incidencia, $\cos\phi_s=\varepsilon\sin^2\theta/(\varepsilon-\sin^2\theta)$; vertical roja:
nulo de HH. Discontinua: $\theta_{pB}$ de cada suelo; banda clara: incidencias de la hélice base.}
\label{fig:C_sigma0_suelos}
\end{figure}

El comportamiento es el mismo en los cuatro suelos (Fig.~\ref{fig:C_sigma0_suelos}):
\begin{itemize}[itemsep=1pt]
  \item \textbf{VV:} en la incidencia media de la hélice (53.4°) el mínimo está en
  $\Delta\phi^\ast=140^{\circ}$ (S0, S1), 136° (S2) y 133° (S3). Está 34--39\,dB por debajo de la mono
  en los suelos con pérdidas (en S0, 74\,dB) y se mantiene a menos de 10\,dB del mínimo en una
  franja de 4--8° de $\Delta\phi$ (0.1° en S0). Al subir la humedad se desplaza poco y se estrecha.
  \item \textbf{HH:} nulo exacto en $\Delta\phi=90^{\circ}$ para todos los suelos e incidencias, porque
  $\alpha_{hh}\propto\cos\phi_s$ no depende de $\varepsilon_r$: es el desacople geométrico entre los
  vectores TE del Tx y del Rx, ortogonales a 90°.
\end{itemize}
La sección siguiente muestra cómo se reflejan estos nulos en el SCR del punto escalar, y
la Sección~\ref{sec:C2} comprueba si sobreviven con blancos reales.

\section{SCR en la imagen}

\begin{figure}[H]
\centering
\includegraphics[width=\textwidth]{C_scr_incidencia_rugosa.pdf}
\caption{SCR en el vóxel del blanco frente a la incidencia central (S1, $s=2$\,cm),
familias F1 y F2, TM y TE. Línea roja: pseudo-Brewster.}
\label{fig:C_scr}
\end{figure}

\begin{figure}[H]
\centering
\includegraphics[width=\textwidth]{C_scr_error_angular.pdf}
\caption{SCR frente al error de incidencia respecto de $\theta_{pB}$ (F2, S1, $s=2$\,cm).}
\label{fig:C_err}
\end{figure}

\begin{table}[H]
\centering\small
\caption{SCR en imagen [dB] en el vóxel del blanco: S1, $s=2$\,cm, superficie congelada (3 realizaciones). \colorbox{mejor}{\textbf{verde}}: mejor canal de la fila; \colorbox{peor}{rojo}: SCR $<0$ (la superficie domina). $\leftarrow$: fila en el pseudo-Brewster. CSV: \texttt{expC/expC\_scr\_imagen.csv}.}
\label{tab:C_scr}
\begin{tabular}{@{}llrrrrrr@{}}
\toprule
Pol. & Familia, $\theta_c$ & Mono & Bi 90° & Bi 135° & Ref1 180° & Mono+bi 90° & Mono+bi 135° \\
\midrule
TM & F1, 35° & \cellcolor{peor}$-$0.5 & 3.3 & 0.8 & \cellcolor{peor}$-$59.6 & 2.8 & \cellcolor{mejor}\textbf{5.8} \\
 & F1, 45° & \cellcolor{peor}$-$0.2 & 1.3 & \cellcolor{mejor}\textbf{13.0} & \cellcolor{peor}$-$58.2 & 4.4 & 5.6 \\
 & F1, 53.4° & 3.1 & 3.1 & \cellcolor{mejor}\textbf{9.8} & \cellcolor{peor}$-$55.3 & 4.3 & 9.5 \\
 & \textbf{F1, $\theta_{pB}$=63.45°} $\leftarrow$ & 8.0 & 6.0 & 7.4 & \cellcolor{peor}$-$41.5 & 9.4 & \cellcolor{mejor}\textbf{12.0} \\
 & F1, 75° & 14.8 & 11.6 & 12.4 & \cellcolor{peor}$-$66.3 & 15.6 & \cellcolor{mejor}\textbf{18.0} \\
\cmidrule(l){2-8}
 & F2, 35° & \cellcolor{peor}$-$4.7 & \cellcolor{peor}$-$1.6 & \cellcolor{peor}$-$13.4 & \cellcolor{peor}$-$61.2 & \cellcolor{mejor}\textbf{$-$0.2} & \cellcolor{peor}$-$8.3 \\
 & F2, 45° & 0.5 & \cellcolor{peor}$-$0.9 & 3.4 & \cellcolor{peor}$-$60.1 & 1.6 & \cellcolor{mejor}\textbf{5.6} \\
 & F2, 53.4° & 0.2 & 3.0 & \cellcolor{mejor}\textbf{6.9} & \cellcolor{peor}$-$57.3 & 2.9 & 4.2 \\
 & \textbf{F2, $\theta_{pB}$=63.45°} $\leftarrow$ & 2.5 & 3.5 & 5.1 & \cellcolor{peor}$-$41.5 & \cellcolor{mejor}\textbf{7.8} & 7.3 \\
 & F2, 75° & 8.4 & 5.7 & 8.3 & \cellcolor{peor}$-$66.0 & \cellcolor{mejor}\textbf{12.6} & 12.1 \\
\midrule
TE & F1, 35° & 1.4 & \cellcolor{mejor}\textbf{14.4} & 1.5 & \cellcolor{peor}$-$64.4 & 6.7 & 4.2 \\
 & F1, 45° & 2.4 & \cellcolor{mejor}\textbf{18.6} & \cellcolor{peor}$-$4.9 & \cellcolor{peor}$-$66.9 & 8.8 & 0.7 \\
 & F1, 53.4° & 6.2 & \cellcolor{mejor}\textbf{23.8} & \cellcolor{peor}$-$2.3 & \cellcolor{peor}$-$69.8 & 11.6 & 4.6 \\
 & \textbf{F1, $\theta_{pB}$=63.45°} $\leftarrow$ & 11.5 & \cellcolor{mejor}\textbf{28.9} & \cellcolor{peor}$-$1.3 & \cellcolor{peor}$-$74.4 & 17.1 & 4.7 \\
 & F1, 75° & 18.5 & \cellcolor{mejor}\textbf{35.0} & 6.9 & \cellcolor{peor}$-$81.3 & 24.6 & 13.6 \\
\cmidrule(l){2-8}
 & F2, 35° & \cellcolor{peor}$-$2.6 & \cellcolor{mejor}\textbf{11.7} & \cellcolor{peor}$-$18.1 & \cellcolor{peor}$-$65.8 & 3.6 & \cellcolor{peor}$-$12.0 \\
 & F2, 45° & 3.5 & \cellcolor{mejor}\textbf{18.5} & \cellcolor{peor}$-$9.1 & \cellcolor{peor}$-$68.6 & 9.0 & \cellcolor{peor}$-$3.3 \\
 & F2, 53.4° & 3.8 & \cellcolor{mejor}\textbf{24.6} & \cellcolor{peor}$-$3.9 & \cellcolor{peor}$-$71.5 & 9.5 & 0.2 \\
 & \textbf{F2, $\theta_{pB}$=63.45°} $\leftarrow$ & 6.7 & \cellcolor{mejor}\textbf{28.7} & \cellcolor{peor}$-$1.9 & \cellcolor{peor}$-$75.5 & 12.9 & 4.2 \\
 & F2, 75° & 13.1 & \cellcolor{mejor}\textbf{36.0} & 4.4 & \cellcolor{peor}$-$81.6 & 19.4 & 10.2 \\
\bottomrule
\end{tabular}
\end{table}


\paragraph{Lectura de la Tabla~\ref{tab:C_scr}.} Cuanto más alto el SCR, mejor.
\begin{enumerate}[itemsep=3pt]
  \item \textbf{El pseudo-Brewster sólo se nota en Ref1, y Ref1 es inservible.} Es la única
  columna con un pico en $\theta_{pB}$: en TM pasa de $-55.3$ a $-41.5$\,dB, y en TE no hay pico
  (control). Aun así queda más de 40\,dB por debajo de cualquier otro canal (toda la columna en rojo).
  Con espectro rectangular la conclusión es la misma ($-34.4$\,dB en $\theta_{pB}$).
  \item \textbf{En los demás canales, más incidencia es mejor, sin pico en $\theta_{pB}$.} El SCR
  crece de forma casi monótona de 35° a 75°, igual en TM que en TE. Por ejemplo, mono F1:
  3.1\,$\to$\,8.0\,$\to$\,14.8\,dB en TM y 6.2\,$\to$\,11.5\,$\to$\,18.5\,dB en TE, en
  53.4°, 63.45° y 75°. La causa es que la difusa cae con la incidencia, por la longitud de
  correlación $\ell$ (factor $W(K)$), no Brewster. El máximo de casi todas las columnas está en
  75°, la mayor incidencia probada.
  \item \textbf{TE supera a TM.} En mono TE gana 2--5\,dB en todas las incidencias,
  al contrario de lo que se esperaría si Brewster (TM) importara: en retrodispersión la difusa
  VV es mayor que la HH.
  \item \textbf{Mejor canal con el punto escalar.} En TE, el bi a 90° es el mejor en las 10 filas
  (12--36\,dB según la incidencia), por el nulo de HH a 90°. \emph{Este resultado no se
  sostiene con blancos reales} (Sección~\ref{sec:C2}). En TM, un canal a 135° (bi o mono+bi) es el mejor en 7 de 10 filas, con 3--13\,dB
  sobre la mono.
\end{enumerate}

\begin{table}[H]
\centering\small
\caption{F1 (hélice) frente a F2 (cono): $\mathrm{SCR}_{F1}-\mathrm{SCR}_{F2}$ en la misma incidencia central, polarización y canal (10 pares por canal: 5 incidencias $\times$ TM/TE, Tabla~\ref{tab:C_scr}). \colorbox{mejor}{verde}: F1 gana.}
\label{tab:C_f1f2}
\begin{tabular}{@{}lrrrc@{}}
\toprule
Canal & media [dB] & mín. [dB] & máx. [dB] & F1 gana \\
\midrule
Mono & \cellcolor{mejor}$+$3.4 & $-$1.0 & 6.4 & 8/10 \\
Bi 90° & \cellcolor{mejor}$+$1.7 & $-$1.0 & 5.9 & 8/10 \\
Bi 135° & \cellcolor{mejor}$+$6.2 & 0.7 & 19.5 & 10/10 \\
Ref1 180° & \cellcolor{mejor}$+$1.2 & $-$0.2 & 2.1 & 9/10 \\
Mono+bi 90° & \cellcolor{mejor}$+$2.6 & $-$0.2 & 5.3 & 9/10 \\
Mono+bi 135° & \cellcolor{mejor}$+$5.9 & 0.0 & 16.2 & 10/10 \\
\midrule
Todos & $+$3.5 & & & 54/60 \\
\bottomrule
\end{tabular}
\end{table}


\paragraph{F1 frente a F2.} La hélice (F1) supera al cono (F2) en 54 de las 60 comparaciones
de igual incidencia, polarización y canal, con 3.5\,dB de SCR de media
(Tabla~\ref{tab:C_f1f2}). El eco del blanco es idéntico en las dos familias (diferencia
$\le0.2$\,dB): lo que cambia es cuánto clutter de la superficie llega al vóxel del blanco. La
superficie está 2\,m por encima del blanco, y F2 tiene peor resolución vertical
($\delta_z$ 2.34 frente a 2.08\,m) y peor ISLR (Tabla~\ref{tab:C_costo}), de modo que filtra
más energía de la superficie hacia el blanco. F2, la órbita que impone Brewster, es por tanto
peor también en SCR, la métrica para la que se propuso.

\begin{table}[H]
\centering\small
\caption{Robustez del SCR en imagen [dB] frente al error de incidencia $\delta$ alrededor de $\theta_{pB}$ (TM, $s=2$\,cm). Pares $a$\,/\,$b$: $\delta<0$\,/\,$\delta>0$. \colorbox{mejor}{verde}: mejor canal de la fila; \colorbox{peor}{rojo}: SCR $<0$.}
\label{tab:C_err_scr}
\begin{tabular}{@{}llrrrr@{}}
\toprule
Suelo, familia & $|\delta|$ & Mono & Bi 135° & Ref1 180° & Mono+bi 135° \\
\midrule
S1, F1 & 0° & 8.0 & 7.4 & \cellcolor{peor}$-$41.5 & \cellcolor{mejor}\textbf{12.0} \\
 & 0.5° & 7.3 / 6.8 & 10.4 / 8.5 & \cellcolor{peor}$-$40.4 / $-$43.4 & \cellcolor{mejor}\textbf{11.9 / 10.5} \\
 & 1° & 6.5 / 10.6 & 8.4 / 9.6 & \cellcolor{peor}$-$40.9 / $-$45.4 & \cellcolor{mejor}\textbf{11.3 / 12.9} \\
 & 2° & 6.6 / 9.6 & 9.4 / 12.3 & \cellcolor{peor}$-$43.9 / $-$49.0 & \cellcolor{mejor}\textbf{11.3 / 15.9} \\
 & 5° & 2.8 / 13.4 & 5.2 / 9.2 & \cellcolor{peor}$-$50.6 / $-$56.4 & \cellcolor{mejor}\textbf{7.0 / 13.6} \\
\midrule
S1, F2 & 0° & 2.5 & 5.1 & \cellcolor{peor}$-$41.5 & \cellcolor{mejor}\textbf{7.3} \\
 & 0.5° & 4.0 / 2.5 & 4.6 / 6.8 & \cellcolor{peor}$-$42.2 / $-$42.4 & \cellcolor{mejor}\textbf{8.3 / 6.4} \\
 & 1° & 4.7 / 5.0 & 3.2 / 7.9 & \cellcolor{peor}$-$43.8 / $-$44.2 & \cellcolor{mejor}\textbf{7.8 / 8.7} \\
 & 2° & 2.5 / 5.4 & 7.2 / 7.9 & \cellcolor{peor}$-$47.1 / $-$48.2 & \cellcolor{mejor}\textbf{8.5 / 9.3} \\
 & 5° & 4.8 / 6.0 & 5.1 / 5.3 & \cellcolor{peor}$-$53.0 / $-$56.2 & \cellcolor{mejor}\textbf{9.6 / 8.7} \\
\midrule
S0, F2 & 0° & 39.0 & 35.9 & \cellcolor{mejor}\textbf{61.2} & 43.1 \\
 & 2° & 37.8 / 41.8 & 38.9 / 40.8 & \cellcolor{peor}$-$11.3 / $-$12.4 & \cellcolor{mejor}\textbf{41.7 / 46.6} \\
\bottomrule
\end{tabular}
\end{table}


\paragraph{Robustez al error de incidencia (Tabla~\ref{tab:C_err_scr}).} En S1, con F2, errores
de hasta $\pm5^{\circ}$ cambian el SCR de mono, bi y fusión en unos $\pm3$\,dB sin tendencia (del
orden de la dispersión entre realizaciones). Con F1 la variación llega a $\pm5$\,dB, pero sigue
la tendencia general con la incidencia ($+\delta$ mejora, $-\delta$ empeora), no un efecto de
Brewster. Sólo Ref1 cae al alejarse de $\theta_{pB}$, hasta 15\,dB con $\pm5^{\circ}$. Es decir,
fuera de Ref1 el SCR es robusto al error angular \emph{porque no depende de Brewster}. En el suelo sin pérdidas
(S0), Ref1 llega a 61\,dB en $\theta_{pB}$ exacto y cae a $-12$\,dB con $\pm2^{\circ}$: es el
mínimo extremadamente estrecho que §2.1 descarta.

\paragraph{Otros factores (una línea cada uno).}
\begin{itemize}[itemsep=1pt]
  \item \textbf{Rugosidad:} el SCR de mono y bi intermedio cae $\approx6$\,dB cada vez que $s$ se
  duplica (difusa $\propto s^2$); el de Ref1 casi no cambia (domina la especular).
  \item \textbf{Suelos húmedos} (S2, S3): SCR negativo en todas las geometrías ($-14$ a $-28$\,dB
  en S2 y $-58$ a $-84$\,dB en S3), porque domina la atenuación del blanco a 2\,m (13.6 y
  27.7\,dB/m por trayecto, frente a 7.7\,dB/m en S1). El pseudo-Brewster no lo remedia.
  \item \textbf{Atitud:} un roll de 0--5° no cambia el SCR en S1 ($\le0.2$\,dB).
\end{itemize}

\section{Costo en imagen de la restricción de incidencia}

\begin{figure}[H]
\centering
\includegraphics[width=\textwidth]{C_costo_imagen.pdf}
\caption{Resolución, volumen y PSLR frente a la incidencia central para F1 (continuo)
y F2 (discontinuo), mono y mono+bi 135°.}
\label{fig:C_costo}
\end{figure}

\begin{table}[H]
\centering\small
\caption{Hélice base (F1 en 53.4°) frente a la restricción de pseudo-Brewster (F2 en
$\theta_{pB}$), con ponderación Hann (\texttt{expC\_psf\_vs\_incidencia.csv},
\texttt{expC\_scr\_imagen.csv}). En cada canal, \colorbox{mejor}{verde} marca la trayectoria que
gana cada métrica. Menor es mejor salvo en contraste y SCR.}
\label{tab:C_costo}
\begin{tabular}{@{}llrrrrrr@{}}
\toprule
Trayectoria & Canal & $\delta_x$ [m] & $\delta_z$ [m] & $V_{-3}$ [l] & ISLR [dB] & Contraste [dB] & SCR [dB] \\
\midrule
F1 base (53.4°) & mono & 0.157 & \cellcolor{mejor}2.08 & 29.0 & \cellcolor{mejor}6.5 & \cellcolor{mejor}19.6 & \cellcolor{mejor}3.1 \\
F2 cono $\theta_{pB}$ & mono & \cellcolor{mejor}0.140 & 2.34 & \cellcolor{mejor}25.0 & 8.3 & 18.4 & 2.5 \\
\midrule
F1 base (53.4°) & mono+bi 135° & 0.213 & \cellcolor{mejor}2.08 & 52.0 & \cellcolor{mejor}2.4 & \cellcolor{mejor}20.4 & \cellcolor{mejor}9.5 \\
F2 cono $\theta_{pB}$ & mono+bi 135° & \cellcolor{mejor}0.191 & 2.34 & \cellcolor{mejor}46.0 & 3.6 & 19.9 & 7.3 \\
\bottomrule
\end{tabular}
\end{table}

\textbf{Costo (pregunta 7).} Imponer la órbita de pseudo-Brewster (F2 en $\theta_{pB}$) sólo gana en
horizontal: $\delta_x$ mejora un 11\,\% (mayor $\sin\theta$) y el volumen del lóbulo un
12--14\,\%. Pierde en todo lo demás: $\delta_z$ empeora un 12.5\,\% (se pierde la diversidad de
elevación de la hélice), el ISLR sube 1.2--1.8\,dB, el contraste baja 0.5--1.2\,dB y el SCR baja
0.6--2.2\,dB. Además obliga a bajar la plataforma a $\approx85$\,m.

\section{Exp.~C2: ¿sobreviven los nulos con blancos reales?}
\label{sec:C2}

\paragraph{Motivación.} El SCR de la Tabla~\ref{tab:C_scr} usa un punto escalar ($S=1$), que no
tiene polarización: sólo la superficie siente los nulos de VV y HH. Un blanco real también
dispersa con polarización, y el mismo desacople geométrico puede anular su eco. El Exp.~C2
separa ambos efectos.

\paragraph{Método.} Tx en la hélice base (53.4°), suelo S1, $s=2$\,cm, $\ell=20$\,cm,
ponderación Hann y 3 realizaciones de la superficie congelada. Rx en la misma hélice con
$\Delta\phi=0^{\circ}{:}15^{\circ}{:}180^{\circ}$ y 140°, TM y TE co-polares. Cuatro blancos en $(0,0,-2)$\,m:
punto escalar, esfera PEC de $r=0.5$\,m, cubo PEC de 1\,m en orientación nominal y el mismo cubo
en 24 orientaciones SO(3). Métricas: SCR en el vóxel del blanco, como en el Exp.~C, y PSCR,
el pico del blanco bajo la superficie dividido por el pico de la superficie en toda la rejilla,
para el canal bi aislado y para la fusión con el mismo presupuesto (N). Script
\texttt{exp\_C2\_nulos.py}; CSV \texttt{expC2/expC2\_scr\_vs\_dphi.csv}. La rejilla y las
realizaciones difieren de las del Exp.~C, de modo que los valores absolutos del punto difieren
hasta 3\,dB de la Tabla~\ref{tab:C_scr}; las comparaciones se hacen dentro de esta tabla.

\begin{figure}[H]
\centering
\includegraphics[width=\textwidth]{C2_nulos_scr.pdf}
\caption{Exp.~C2 frente a $\Delta\phi$ (hélice base, S1, $s=2$\,cm). (a, d) Potencia en la
imagen del canal bi respecto de la mono: superficie en el vóxel del blanco (rojo) y eco de cada
blanco. (b, e) Ganancia de SCR en el vóxel respecto de la mono. (c, f) Ganancia de PSCR (pico
del blanco / pico de la superficie). Continua: bi aislado; discontinua: fusión mono+bi (N).
Líneas grises en 90° (nulo HH) y 140° (nulo VV).}
\label{fig:C2}
\end{figure}

\begin{table}[H]
\centering\small
\caption{Exp.~C2: SCR en el vóxel del blanco [dB]. ``Superficie'': potencia de la superficie
en el vóxel respecto de la mono. \colorbox{mejor}{verde}: mejora de más de 3\,dB sobre la mono
de la misma polarización; \colorbox{peor}{rojo}: empeora más de 3\,dB.}
\label{tab:C2}
\begin{tabular}{@{}llrrrrr@{}}
\toprule
Pol. & Canal & Superficie & Punto & Esfera & Cubo nominal & Cubo SO(3) \\
\midrule
TM & mono & 0.0 & 0.7 & $-0.9$ & $-9.5$ & $-12.0$ \\
 & bi 90° & $-2.0$ & 2.6 & \cellcolor{peor}$-13.3$ & \cellcolor{peor}$-25.3$ & \cellcolor{peor}$-27.9$ \\
 & bi 120° & $-4.6$ & \cellcolor{mejor}5.2 & \cellcolor{peor}$-6.1$ & \cellcolor{peor}$-19.1$ & \cellcolor{mejor}$-6.2$ \\
 & bi 135° & $-12.5$ & \cellcolor{mejor}13.1 & \cellcolor{mejor}7.5 & \cellcolor{mejor}14.7 & \cellcolor{mejor}6.2 \\
 & bi 140° & $-13.6$ & \cellcolor{mejor}14.2 & \cellcolor{mejor}\textbf{9.8} & \cellcolor{mejor}\textbf{20.1} & \cellcolor{mejor}\textbf{8.3} \\
 & bi 150° & $-1.2$ & 1.8 & $-0.9$ & \cellcolor{mejor}13.7 & \cellcolor{mejor}$-2.7$ \\
 & mono+bi 140° & $-5.7$ & \cellcolor{mejor}6.4 & $-0.9$ & \cellcolor{mejor}4.7 & \cellcolor{mejor}$-4.2$ \\
\midrule
TE & mono & 0.0 & 3.9 & 2.3 & $-6.4$ & $-8.8$ \\
 & bi 90° & $-17.8$ & \cellcolor{mejor}\textbf{21.7} & \cellcolor{peor}$-4.8$ & $-6.2$ & $-9.0$ \\
 & mono+bi 90° & $-5.5$ & \cellcolor{mejor}9.4 & 2.2 & $-7.4$ & $-9.3$ \\
\bottomrule
\end{tabular}
\end{table}

\paragraph{Resultados e interpretación.}
\begin{enumerate}[itemsep=3pt]
  \item \textbf{El nulo HH en TE a 90° es un artefacto del punto escalar.} La superficie cae
  17.8\,dB, pero el eco de la esfera cae 25\,dB y el del cubo SO(3) 18\,dB. El SCR de la esfera
  empeora 7.1\,dB respecto de la mono y el del cubo SO(3) no cambia. Es el mismo mecanismo que el
  de la Sección~\ref{sec:D3}: a 90° los vectores TE del Tx y del Rx son ortogonales, y eso
  desacopla tanto la superficie como el blanco. Sólo el punto, sin polarización, gana 17.8\,dB.
  \item \textbf{El nulo VV en TM a 135°--140° sí sobrevive.} La superficie cae 12.5 a 13.6\,dB,
  mientras que el eco de los blancos apenas cae (esfera, $-4.1$ y $-2.9$\,dB) o incluso sube
  (cubo SO(3), $+5.8$ y $+6.7$\,dB, por la especular de sus caras). El SCR del canal bi aislado
  mejora respecto de la mono 8.4 a 10.7\,dB en la esfera, 24 a 30\,dB en el cubo nominal y 18 a
  20\,dB en el cubo con orientación aleatoria. En PSCR la mejora es menor pero del mismo signo:
  $+2.5$ a $+3.7$\,dB en la esfera y $+12$ a $+14$\,dB en el cubo SO(3).
  \item \textbf{La ventana es estrecha.} A 120° y a 150° la esfera ya no mejora ($-6.1$ y
  $-0.9$\,dB) y desde 150° la especular de nadir domina el PSCR. El Rx debe mantenerse a
  $\approx\pm5^{\circ}$ de 140°, coherente con la anchura de 4--8° del nulo VV
  (Fig.~\ref{fig:C_sigma0_suelos}).
  \item \textbf{La fusión de igual peso desperdicia el canal limpio.} La imagen fusionada es la
  media de las dos, y el clutter de la mono la domina: su superficie sólo baja 5.7\,dB, casi el
  límite de 6\,dB de una media de dos canales. A 140° la fusión deja la esfera igual que la mono y
  queda 10 a 15\,dB por debajo del bi aislado. Para aprovechar el nulo hay que usar el canal bi
  solo o una fusión ponderada por el clutter de cada canal.
  \item \textbf{El precio es la resolución.} A 135° el canal bi tiene $\delta_x=0.41$\,m frente a
  0.157\,m en la mono (Tabla~\ref{tab:expB}), con la misma $\delta_z$. Es un modo de detección
  en escenarios limitados por clutter, no un modo de imagen fina.
\end{enumerate}

\conclexp{C2}{Con blancos reales, la única geometría que suprime el clutter de forma útil es
TM con el Rx a 135°--140° usado como canal aislado: el SCR mejora 8 a 20\,dB sobre la mono, a
cambio de una resolución horizontal 2.6 veces peor. El nulo HH en TE a 90° es un artefacto del
punto escalar.}

\section{Conclusiones del Exp.~C: ¿qué es mejor?}

\begin{table}[H]
\centering\small
\caption{Respuesta directa a las preguntas del Exp.~C (S1, $\varepsilon_r=3.99-j0.40$, $s=2$\,cm).
\colorbox{mejor}{verde}: la opción recomendada.}
\label{tab:C_concl}
\begin{tabular}{@{}L{3.6cm}L{3.2cm}L{3.4cm}L{4.4cm}@{}}
\toprule
\textbf{Pregunta} & \textbf{Mejor opción} & \textbf{Margen} & \textbf{Por qué} \\
\midrule
¿Volar restringido al pseudo-Brewster? & \cellcolor{mejor}No & F2 en $\theta_{pB}$ pierde 0.6--2.2\,dB de SCR (mono y mono+bi 135°) frente a la hélice base &
Sólo suprime la especular, que sólo pesa en Ref1 \\
¿Hélice (F1) o cono (F2)? & \cellcolor{mejor}F1 (hélice) & $+3.5$\,dB de SCR medio; gana 54/60 &
Mejor $\delta_z$ e ISLR: menos superficie en el vóxel del blanco \\
¿Qué incidencia? & \cellcolor{mejor}La mayor viable (75° fue la mejor probada) & Mono TM: $+11.7$\,dB de 53.4° a 75° &
La difusa cae con $\theta$ ($\ell$); igual en TM y TE \\
¿Qué canal en TM? & \cellcolor{mejor}Bi aislado a 135°--140° & 8--20\,dB de SCR sobre la mono con esfera y cubo SO(3) (Exp.~C2) & Nulo VV biestático; el eco del blanco no cae (SPM; hay que validarlo) \\
¿TM o TE? & \cellcolor{mejor}TM & TE a 90° sólo gana con el punto escalar; con la esfera pierde 7\,dB & El desacople a 90° anula también el eco del blanco \\
¿Fusión de igual peso? & \cellcolor{peor}No para clutter & Queda 10--15\,dB por debajo del bi aislado a 140° & El clutter de la mono domina la media \\
¿Ref1 (180°)? & \cellcolor{peor}Nunca & $-41$ a $-81$\,dB & La especular cae sobre el blanco \\
¿Es robusto al error angular? & Sí, salvo Ref1 & F2: $\pm3$\,dB con $\pm5^{\circ}$ & Porque no depende de Brewster \\
\bottomrule
\end{tabular}
\end{table}

\paragraph{En una frase.} El pseudo-Brewster no es un criterio útil de trayectoria: la órbita
que lo impone (F2) es peor que la hélice en SCR, resolución vertical y lóbulos, y sólo mejora
Ref1, que sigue inservible. Lo que sí sube el SCR es la geometría, no Brewster: mayor incidencia y, sobre todo,
el Rx a 135°--140° en TM usado como canal aislado, que mejora 8 a 20\,dB el SCR de blancos reales.

\paragraph{Cautelas.} (i) La incidencia de 75° se evaluó sin restricciones operativas: el
Tx volaría a $\approx49$\,m de altura, con más alcance y otra cobertura del haz. (ii) El nulo
de VV es de SPM de primer orden y debe verificarse con un modelo de orden superior (IEM/AIEM),
con polarización cruzada y con medidas. (iii) El Exp.~C2 usa 3 realizaciones de la superficie y
24 orientaciones del cubo; las ganancias de 8 a 20\,dB superan con margen la dispersión de
$\pm3$\,dB.

% ══════════════════════════════════════════════════════════════════════════
\chapter{Experimento D: biblioteca de blancos}
\label{cap:resD}

\paragraph{Hipótesis.} Las mejoras de lóbulos y contraste del Exp.~B (punto
isótropo) se mantienen para blancos extensos con RCS anisótropa.

\paragraph{Método.} PEC por PO biestática, centro en $(0,0,-2)$\,m, suelo S1,
orientación nominal (P0: esfera $r=0.5$\,m, cubo 1\,m, caja $1\times0.5\times0.3$\,m,
cilindro vertical $r=0.3$, $L=1$\,m; P1: tubo horizontal $r=0.15$, $L=2$\,m,
tubo inclinado 30° sobre la horizontal y 30° de acimut, cilindro inclinado 45°;
P2: escena con cubo, tubo y esfera). Rejilla de imagen $3\times3\times4$\,m. CSV:
\texttt{expD\_objetos.csv}, \texttt{expD\_escenaP2.csv},
\texttt{expD2\_rejilla\_densa.csv}. Para objetos extensos las ``métricas de PSF'' se
evalúan sobre la respuesta del objeto (pico, lóbulo, contraste) y se añade la
fracción de energía dentro de la región del objeto (caja envolvente $+0.3$\,m).

\begin{figure}[H]
\centering
\includegraphics[width=0.92\textwidth]{D_mip_objetos.pdf}
\caption{Proyección de máxima intensidad (en $y$) de cada blanco: mono, bi y
fusión para $\Delta\phi=90^{\circ}$ y 135°. Escala común por blanco (pico mono $=0$\,dB, 30\,dB).}
\label{fig:D_mip}
\end{figure}

\begin{figure}[H]
\centering
\includegraphics[width=\textwidth]{D_fusion_vs_mono.pdf}
\caption{Métricas de la fusión con el mismo presupuesto (N) respecto de la
imagen monoestática, por blanco y por $\Delta\phi$.}
\label{fig:D_fus}
\end{figure}

\begin{table}[H]
\centering\small
\caption{Tabla comparativa mono vs.\ bi vs.\ fusión para los casos P0/P1 en
orientación nominal (diferencias fusión(N) $-$ mono, dB salvo energía en puntos
porcentuales). ``Bi pico'': pico del canal bi respecto del mono.}
\label{tab:P0}
\begin{tabular}{@{}lr|r|rrrrrr@{}}
\toprule
Blanco & $\Delta\phi$ & Bi pico & $\Delta V$ & $\Delta$PSLR & $\Delta$ISLR & $\Delta$Contr. & $\Delta$SNR & $\Delta E$ \\
\midrule
punto & 90 & $+0.2$ & $+1.1$ & $\mathbf{-2.0}$ & $\mathbf{-4.0}$ & $\mathbf{+2.6}$ & $+0.1$ & $+4.1$ \\
 & 135 & $+0.3$ & $+2.4$ & $-2.4$ & $-2.7$ & $-0.4$ & $+0.1$ & $+0.8$ \\
esfera & 90 & $-14.9$ & $+0.2$ & $+0.7$ & $-1.5$ & $+1.3$ & $-4.8$ & $-0.2$ \\
 & 135 & $-3.5$ & $+0.9$ & $+2.6$ & $+3.0$ & $-4.3$ & $-5.8$ & $-8.5$ \\
cubo & 90 & $-19.3$ & $+0.1$ & $+0.3$ & $+0.9$ & $-0.8$ & $-6.7$ & $+1.0$ \\
 & 135 & $+12.6$ & $+11.7$ & $+3.1$ & $+2.5$ & $-4.8$ & $+5.1$ & $-6.3$ \\
caja & 90 & $-10.5$ & $+1.2$ & $+0.5$ & $-0.5$ & $-0.7$ & $-6.3$ & $+1.5$ \\
 & 135 & $+17.4$ & $+8.5$ & $-2.8$ & $-7.6$ & $-1.1$ & $+11.7$ & $-8.2$ \\
cilindro vertical & 90 & $-16.9$ & $+0.1$ & $-1.0$ & $+0.2$ & $-0.5$ & $-6.3$ & $+0.4$ \\
 & 135 & $+16.7$ & $+5.7$ & $-2.0$ & $-3.8$ & $-3.3$ & $+11.9$ & $-11.7$ \\
tubo horizontal & 90 & $-19.6$ & $+0.2$ & $0.0$ & $-0.8$ & $+0.6$ & $-5.4$ & $+2.1$ \\
 & 135 & $+6.2$ & $+4.6$ & $-0.2$ & $-8.3$ & $+0.7$ & $+1.7$ & $-4.5$ \\
tubo inclinado & 90 & $-19.9$ & $+0.1$ & $+0.2$ & $-0.2$ & $+0.1$ & $-5.8$ & $+0.8$ \\
 & 135 & $-5.7$ & $+3.7$ & $+0.4$ & $+1.0$ & $-3.2$ & $-6.8$ & $-4.2$ \\
cilindro inclinado & 90 & $-16.9$ & $+0.2$ & $-0.1$ & $-0.9$ & $+0.7$ & $-5.1$ & $+1.3$ \\
 & 135 & $-4.2$ & $+3.1$ & $+3.2$ & $+0.6$ & $-3.3$ & $-5.7$ & $-5.1$ \\
\midrule
todos & 180 & \multicolumn{7}{l}{fusión siempre peor: $\Delta V=+6$ a $+28$\,dB, $\Delta$contraste $-8$ a $-17$\,dB} \\
\bottomrule
\end{tabular}
\end{table}

\paragraph{Resultados e interpretación.}
\begin{enumerate}[itemsep=3pt]
  \item \textbf{La ganancia de lóbulos del Exp.~B es propia del dispersor
  isótropo.} Sólo el punto (y, parcialmente, la esfera a 90°) reproduce la
  mejora de PSLR/ISLR/contraste. En los objetos con caras planas, la respuesta
  está dominada por las reflexiones especulares de las caras, que ya fijan
  el ``PSLR'' de la imagen mono en $-0.5$ a $-4.7$\,dB y no dependen del
  relleno del dominio $k$.
  \item \textbf{Desbalance de canales.} A 90° el eco bi de los objetos es
  $10$--$20$\,dB más débil que el mono, incluso en la esfera. La causa es un desajuste
  de polarización del Rx TM (Sección~\ref{sec:D3}). La fusión de igual peso
  resulta prácticamente la imagen mono con $5$--$7$\,dB menos de SNR (con el
  mismo presupuesto N): \emph{la fusión de igual peso diluye la señal}.
  \item A 135° y 180° ocurre lo contrario con cubo, caja, cilindro vertical y
  tubo horizontal: la bisectriz casi vertical mantiene la especular de la cara
  superior (efecto ya visto el 28/09) y el canal bi domina ($+6$ a $+30$\,dB);
  la fusión hereda entonces el lóbulo ancho del canal bi ($\Delta V$ hasta
  $+12$\,dB a 135° y $+28$\,dB a 180°). Más eco no es mejor imagen
  (§11: ``Não assumir que maior RCS biestática implica melhor resolução'').
  \item \textbf{Escena P2 (enmascaramiento).} En mono la esfera domina y el tubo
  queda 18.1\,dB por debajo, bajo los lóbulos de los otros objetos
  (relación objeto--lóbulo mínima $-17.5$\,dB). Con la fusión a 135° el rango
  dinámico entre objetos se reduce de 18.1 a 7.9\,dB y el peor objeto mejora
  10.1\,dB respecto de los lóbulos ($-7.4$\,dB): la diversidad de aspecto del
  canal bi reduce el enmascaramiento (Fig.~\ref{fig:D_P2}). Es el beneficio de
  complementariedad más claro observado con objetos reales. A 90° la fusión no
  cambia nada ($-17.7$\,dB).
\end{enumerate}

\begin{figure}[H]
\centering
\includegraphics[width=\textwidth]{D_escenaP2.pdf}
\caption{Escena P2 (cubo, tubo inclinado en acimut y esfera de distinta RCS):
MIP en $z$, misma escala.}
\label{fig:D_P2}
\end{figure}

\section{Separabilidad de blancos próximos}
\label{sec:D2}
Con blancos separados 1\,m, todos los canales salvo Ref1 los resuelven todos. Por eso la
separabilidad se midió con blancos más próximos en un plano horizontal: 16 puntos en
$z=-2$\,m (S0) con separación $s=0.15$--$0.30$\,m ($\delta_x^{mono}=0.157$\,m), fases
aleatorias fijas por punto y 3 realizaciones (\texttt{exp\_D2\_rejilla\_densa.py}).

\begin{table}[H]
\centering\small
\caption{Fracción de blancos resueltos (caída $\ge3$\,dB hacia cada vecino),
media de 3 realizaciones (\texttt{expD2\_separabilidad\_resumen.csv}).}
\label{tab:D2}
\begin{tabular}{@{}lrrrr@{}}
\toprule
Canal & $s=0.15$\,m & 0.20\,m & 0.25\,m & 0.30\,m \\
\midrule
Mono & 8\,\% & 8\,\% & 29\,\% & \textbf{73\,\%} \\
Mono+bi 45° & 6\,\% & 2\,\% & 25\,\% & 63\,\% \\
Mono+bi 90° & 6\,\% & 0\,\% & 10\,\% & 27\,\% \\
Mono+bi 120° & 2\,\% & 2\,\% & 10\,\% & 13\,\% \\
Mono+bi 135° & 2\,\% & 0\,\% & 13\,\% & 15\,\% \\
Bi 90° aislado & 0\,\% & 0\,\% & 4\,\% & 17\,\% \\
\bottomrule
\end{tabular}
\end{table}

\begin{figure}[H]
\centering
\includegraphics[width=\textwidth]{D2_separabilidad.pdf}
\caption{Separabilidad horizontal vs.\ separación e imágenes con $s=0.20$\,m.}
\label{fig:D2}
\end{figure}

\textbf{Resultado negativo:} la separación de blancos próximos la gobierna el
ancho del lóbulo principal, y la fusión lo ensancha. La imagen mono separa
mejor que cualquier fusión; la fusión a 90°--135° resuelve entre 2.7 y 5.6
veces menos blancos a $s=0.30$\,m. La menor sidelobe de la fusión no compensa
el ensanchamiento cuando los blancos están a $\approx2\delta_x$.


\section{Rx con polarización alineada}
\label{sec:D3}

Con un Rx TM co-polar, el pico del canal bi de los objetos a $\Delta\phi=90^{\circ}$ cae
10--20\,dB respecto de la mono (Tabla~\ref{tab:P0}); el Monte Carlo del Exp.~E confirma
$\approx20$\,dB en potencia media para cualquier orientación. La causa es geométrica: dentro del suelo
los rayos forman $\approx23^{\circ}$ con la vertical, de modo que el vector TM es casi horizontal y
apunta radialmente hacia cada UAV; a $\Delta\phi=90^{\circ}$ el TM del Rx es casi ortogonal al
campo que dispersan las caras del objeto (acoplamiento $\propto\cos\Delta\phi$). El punto
isótropo escalar no ve este efecto. Se evaluó entonces un Rx \emph{alineado}: el Tx
transmite TM y el Rx, de doble polarización, combina sus canales para recibir la
componente del campo del Tx transversal a la dirección de recepción (en mono coincide
con TM--TM); la transmitancia de subida pondera TM y TE según esa proyección
(\texttt{exp\_D3\_rx\_alineado.py}, \texttt{expD3\_comparacion\_TM\_vs\_alineado.csv}).

\begin{figure}[H]
\centering
\includegraphics[width=\textwidth]{D3_rx_alineado.pdf}
\caption{Pico del canal bi respecto del mono con Rx TM co-polar y con Rx alineado (orientación nominal).}
\label{fig:D3}
\end{figure}

\begin{table}[H]
\centering\small
\caption{Rx TM frente a Rx alineado a $\Delta\phi=90^{\circ}$ (fusión con el mismo presupuesto N, diferencias respecto de mono).}
\label{tab:D3}
\begin{tabular}{@{}lrr|rr|rr@{}}
\toprule
 & \multicolumn{2}{c|}{Pico bi [dB]} & \multicolumn{2}{c|}{$\Delta$SNR fus [dB]} & \multicolumn{2}{c}{$\Delta$contraste fus [dB]} \\
Blanco & TM & alin. & TM & alin. & TM & alin. \\
\midrule
esfera & $-14.9$ & $-1.9$ & $-4.8$ & $-1.5$ & $+1.3$ & $+2.0$ \\
cubo & $-19.3$ & $+8.4$ & $-6.7$ & $+0.2$ & $-0.8$ & $-3.7$ \\
caja & $-10.5$ & $+9.3$ & $-6.3$ & $+3.6$ & $-0.7$ & $-0.5$ \\
cilindro vertical & $-16.9$ & $+9.8$ & $-6.3$ & $+0.5$ & $-0.5$ & $-5.5$ \\
tubo horizontal & $-19.6$ & $+1.8$ & $-5.4$ & $+0.7$ & $+0.6$ & $+2.3$ \\
tubo inclinado & $-19.9$ & $-2.7$ & $-5.8$ & $-1.2$ & $+0.1$ & $+3.0$ \\
cilindro inclinado & $-16.9$ & $+0.5$ & $-5.1$ & $-0.1$ & $+0.7$ & $+2.1$ \\
\bottomrule
\end{tabular}
\end{table}

Con el Rx alineado el canal bi a 90° recupera entre 13 y 28\,dB y la fusión deja de diluir la
SNR ($-1.5$ a $+3.6$\,dB frente a $-4.8$ a $-6.7$\,dB). El contraste de la fusión mejora
2--3\,dB en la esfera, los tubos y el cilindro inclinado; en los objetos con una cara
horizontal (cubo, caja, cilindro vertical) el especular biestático pasa a dominar y el
contraste empeora. \textbf{Consecuencia de diseño:} el Rx remoto debe ser de doble polarización
(o girar su antena con $\Delta\phi$); un Rx TM co-polar anula la complementariedad justo en la
región de $\Delta\phi$ que el dominio $k$ recomienda.

\conclexp{D}{Con objetos reales la fusión sólo ayuda cuando el canal bi llega con potencia
comparable a la mono: reduce 10.1\,dB el enmascaramiento en la escena P2 a 135°, pero con un Rx
TM co-polar a 90° el bi pierde 10--20\,dB por desajuste de polarización (un Rx alineado recupera
13--28\,dB), y la mono separa mejor blancos próximos (73\,\% frente a 27\,\% a 0.30\,m).}

% ══════════════════════════════════════════════════════════════════════════
\chapter{Experimento E: orientación y robustez}
\label{cap:resE}

\paragraph{Hipótesis.} Las mejoras de la fusión observadas con el punto y con los
objetos en orientación nominal se mantienen cuando la orientación del objeto es
desconocida.

\paragraph{Método (§5.6).} Monte Carlo \emph{por adquisición}: en cada repetición se
sortea una rotación uniforme en SO(3) (cuaternión gaussiano normalizado), se aplica
al objeto durante toda la apertura y es la misma en los canales mono y bi (diseño
pareado). Cubo (1000 orientaciones), caja, cilindro y tubo (200 cada uno),
$\Delta\phi=45^{\circ},90^{\circ},135^{\circ},180^{\circ}$, 800 pulsos, rejilla de $2.4\times2.4\times3$\,m. Se
corrigió así la práctica del 28/09 de sortear una orientación por posición. Se
corrieron dos receptores: TM co-polar (Exp.~E) y alineado de doble polarización
(Exp.~E2, $\Delta\phi=45^{\circ},90^{\circ},135^{\circ}$, 200 orientaciones por objeto). CSV:
\texttt{expE/expE\_mc\_<objeto>.csv} (cada repetición, con su matriz de rotación),
\texttt{expE/expE\_resumen.csv} y \texttt{expE2/expE\_resumen.csv} (medianas, IC\,95\,\%
bootstrap y probabilidad de mejora).

% generado por simulations/analise_1610/tabla_E_latex.py
\begin{table}[H]
\centering\small
\caption{Monte Carlo SO(3) por adquisici\'on: mediana de (fusi\'on N $-$ mono) y probabilidad de mejora P$_{mej}$ (Rx TM co-polar). $\Delta V$ en dB; loc.\ en m. N = 1000 (cubo), 200.}
\label{tab:E_TM}
\resizebox{\textwidth}{!}{\begin{tabular}{@{}lrr|rrrrrrrrrr@{}}
\toprule
Objeto & $\Delta\phi$ & $\bar P_{bi}/\bar P_{mono}$ & \multicolumn{2}{c}{$\Delta$SNR} & \multicolumn{2}{c}{$\Delta$Contr.} & \multicolumn{2}{c}{$\Delta$PSLR} & \multicolumn{2}{c}{$\Delta V$} & \multicolumn{2}{c}{$\Delta$loc} \\
 & & [dB] & med. & P$_{mej}$ & med. & P$_{mej}$ & med. & P$_{mej}$ & med. & P$_{mej}$ & med. & P$_{mej}$ \\
\midrule
cubo & 45° & -2.6 & -1.4 & 7\,\% & +0.3 & 81\,\% & +0.00 & 49\,\% & +0.1 & 38\,\% & +0.004 & 36\,\% \\
cubo & 90° & -20.2 & -5.7 & 0\,\% & +0.3 & 78\,\% & +0.02 & 47\,\% & +0.1 & 39\,\% & -0.001 & 59\,\% \\
cubo & 135° & -3.8 & -4.4 & 13\,\% & -0.0 & 48\,\% & +0.07 & 42\,\% & +0.9 & 43\,\% & +0.049 & 29\,\% \\
cubo & 180° & -0.0 & -2.9 & 31\,\% & -2.8 & 10\,\% & +0.11 & 34\,\% & -0.4 & 53\,\% & +0.014 & 42\,\% \\
\midrule
caja & 45° & -2.4 & -1.2 & 2\,\% & +0.4 & 98\,\% & -0.06 & 71\,\% & +0.1 & 33\,\% & +0.001 & 42\,\% \\
caja & 90° & -19.8 & -5.4 & 0\,\% & +0.6 & 91\,\% & -0.00 & 52\,\% & +0.1 & 40\,\% & -0.000 & 51\,\% \\
caja & 135° & -5.2 & -4.8 & 10\,\% & -0.5 & 37\,\% & +0.06 & 46\,\% & +0.6 & 44\,\% & +0.049 & 34\,\% \\
caja & 180° & -2.9 & -2.4 & 28\,\% & -5.2 & 3\,\% & +0.19 & 34\,\% & +0.7 & 43\,\% & +0.045 & 36\,\% \\
\midrule
cilindro & 45° & -2.6 & -1.1 & 0\,\% & +0.5 & 100\,\% & -0.00 & 50\,\% & +0.2 & 25\,\% & +0.000 & 46\,\% \\
cilindro & 90° & -20.3 & -5.5 & 0\,\% & +0.6 & 99\,\% & -0.06 & 66\,\% & +0.1 & 30\,\% & -0.001 & 60\,\% \\
cilindro & 135° & -3.1 & -5.1 & 7\,\% & -1.3 & 21\,\% & +0.48 & 26\,\% & +3.5 & 22\,\% & +0.002 & 49\,\% \\
cilindro & 180° & +0.9 & -1.6 & 28\,\% & -7.2 & 0\,\% & +0.74 & 19\,\% & +4.4 & 26\,\% & +0.061 & 14\,\% \\
\midrule
tubo & 45° & -2.4 & -1.1 & 3\,\% & +0.4 & 96\,\% & -0.00 & 50\,\% & +0.2 & 11\,\% & -0.002 & 74\,\% \\
tubo & 90° & -19.5 & -5.4 & 0\,\% & +0.6 & 94\,\% & +0.02 & 46\,\% & +0.1 & 23\,\% & -0.000 & 57\,\% \\
tubo & 135° & -3.7 & -5.6 & 16\,\% & -2.0 & 28\,\% & +0.42 & 27\,\% & +2.3 & 24\,\% & +0.195 & 40\,\% \\
tubo & 180° & -0.1 & -2.6 & 26\,\% & -7.9 & 0\,\% & +0.51 & 22\,\% & +4.5 & 31\,\% & +0.020 & 42\,\% \\
\bottomrule
\end{tabular}}
\end{table}
\begin{table}[H]
\centering\small
\caption{Monte Carlo SO(3) por adquisici\'on: mediana de (fusi\'on N $-$ mono) y probabilidad de mejora P$_{mej}$ (Rx alineado, doble polarizaci\'on). $\Delta V$ en dB; loc.\ en m. N = 200.}
\label{tab:E_AL}
\resizebox{\textwidth}{!}{\begin{tabular}{@{}lrr|rrrrrrrrrr@{}}
\toprule
Objeto & $\Delta\phi$ & $\bar P_{bi}/\bar P_{mono}$ & \multicolumn{2}{c}{$\Delta$SNR} & \multicolumn{2}{c}{$\Delta$Contr.} & \multicolumn{2}{c}{$\Delta$PSLR} & \multicolumn{2}{c}{$\Delta V$} & \multicolumn{2}{c}{$\Delta$loc} \\
 & & [dB] & med. & P$_{mej}$ & med. & P$_{mej}$ & med. & P$_{mej}$ & med. & P$_{mej}$ & med. & P$_{mej}$ \\
\midrule
cubo & 45° & -0.5 & -0.5 & 25\,\% & +0.3 & 80\,\% & +0.02 & 46\,\% & +0.2 & 35\,\% & +0.003 & 41\,\% \\
cubo & 90° & -1.4 & -2.1 & 24\,\% & +1.8 & 90\,\% & -0.05 & 53\,\% & +0.8 & 33\,\% & -0.009 & 61\,\% \\
cubo & 135° & -0.5 & -2.9 & 30\,\% & +1.2 & 78\,\% & -0.06 & 59\,\% & +0.8 & 40\,\% & -0.001 & 52\,\% \\
\midrule
caja & 45° & -0.6 & -0.4 & 22\,\% & +0.5 & 94\,\% & -0.04 & 62\,\% & +0.1 & 28\,\% & +0.001 & 40\,\% \\
caja & 90° & -1.6 & -1.5 & 20\,\% & +2.6 & 98\,\% & -0.28 & 70\,\% & +0.6 & 30\,\% & -0.001 & 54\,\% \\
caja & 135° & -1.0 & -2.7 & 21\,\% & +1.0 & 84\,\% & -0.36 & 70\,\% & +0.7 & 33\,\% & +0.002 & 49\,\% \\
\midrule
cilindro & 45° & -0.5 & -0.2 & 29\,\% & +0.5 & 100\,\% & -0.02 & 57\,\% & +0.2 & 24\,\% & -0.000 & 50\,\% \\
cilindro & 90° & -1.6 & -0.9 & 30\,\% & +2.7 & 96\,\% & -0.39 & 66\,\% & +0.8 & 23\,\% & -0.005 & 66\,\% \\
cilindro & 135° & -1.6 & -1.8 & 26\,\% & +0.9 & 80\,\% & -0.59 & 76\,\% & +1.4 & 24\,\% & -0.000 & 50\,\% \\
\midrule
tubo & 45° & -0.5 & -0.3 & 36\,\% & +0.5 & 98\,\% & -0.04 & 64\,\% & +0.3 & 15\,\% & -0.001 & 59\,\% \\
tubo & 90° & -1.0 & -1.1 & 39\,\% & +2.6 & 98\,\% & -0.08 & 65\,\% & +1.0 & 10\,\% & -0.001 & 56\,\% \\
tubo & 135° & -0.0 & -1.6 & 35\,\% & +1.0 & 92\,\% & -0.07 & 58\,\% & +1.5 & 18\,\% & +0.001 & 42\,\% \\
\bottomrule
\end{tabular}}
\end{table}


\begin{figure}[H]
\centering
\includegraphics[width=\textwidth]{E_montecarlo.pdf}
\caption{Monte Carlo SO(3) con Rx TM co-polar: distribución de (fusión N $-$ mono) por
objeto y separación; porcentajes = probabilidad de mejora.}
\label{fig:E_mc}
\end{figure}

\begin{figure}[H]
\centering
\includegraphics[width=\textwidth]{E2_montecarlo_rx_alineado.pdf}
\caption{Ídem con Rx de polarización alineada (doble polarización).}
\label{fig:E2_mc}
\end{figure}

\paragraph{Resultados e interpretación.}
\begin{enumerate}[itemsep=3pt]
  \item \textbf{Potencia (control).} Con orientación aleatoria, la potencia media del canal
  bi a 180° es la misma que la mono ($-2.9$ a $+0.9$\,dB según el objeto), como concluyó el
  28/09. Con Rx TM co-polar cae $\approx20$\,dB a 90° y 2.4--5.2\,dB a 45° y 135°
  ($\approx\cos^2\Delta\phi$): es el desajuste geométrico de polarización. Con el Rx alineado
  queda entre 0 y $-1.6$\,dB de la mono en todas las separaciones.
  \item \textbf{Rx TM co-polar:} a 45°--90° la fusión gana un contraste pequeño pero
  sistemático ($+0.3$ a $+0.6$\,dB, mejora en el 78--100\,\% de las adquisiciones) y pierde
  1.1--5.7\,dB de SNR (con el mismo presupuesto N; la fusión 2N recupera 3\,dB). PSLR y
  volumen no cambian (medianas $\approx0$). A 135°--180° la fusión empeora el contraste
  (hasta $-7.9$\,dB) y el volumen.
  \item \textbf{Rx alineado:} a 90° el contraste mejora $+1.8$ a $+2.7$\,dB con
  probabilidad del 90--98\,\% en los cuatro objetos; el PSLR mejora ligeramente
  (mediana $-0.05$ a $-0.39$\,dB, 53--70\,\%), la pérdida de SNR baja a $-0.9$ a
  $-2.1$\,dB (N) y la energía en la región del objeto sube (64--86\,\%). El precio es un
  lóbulo $+0.6$ a $+1.0$\,dB más voluminoso. A 135° la ganancia de contraste baja a
  $+0.9$ a $+1.2$\,dB.
  \item \textbf{La mejora se mantiene con la orientación sólo si el Rx está alineado
  en polarización y a 45°--90°.} La resolución y la localización no mejoran en ningún caso
  (medianas $\le0.01$\,m de diferencia).
\end{enumerate}

\conclexp{E}{Con orientación aleatoria del objeto, la fusión mejora el contraste de forma
sistemática sólo con un Rx de polarización alineada a 90° ($+1.8$ a $+2.7$\,dB en el 90--98\,\% de
las orientaciones); con un Rx TM co-polar la ganancia es marginal ($+0.3$ a $+0.6$\,dB) y cuesta
1--6\,dB de SNR, y en ningún caso mejoran la resolución ni la localización.}

% ══════════════════════════════════════════════════════════════════════════
\chapter{Experimento F: trayectoria del Rx para imagen}
\label{cap:resF}

\paragraph{Hipótesis.} Partiendo de la región prometedora del Exp.~B
($\Delta\phi\approx90^{\circ}$--$120^{\circ}$), variar la separación a lo largo de la
adquisición rellena el disco horizontal del dominio $k$ con más anillos y
mejora la concentración de energía de la imagen fusionada, con un costo de
vuelo bajo. El radio del Rx (a la misma altura) podría aportar diversidad
vertical.

\paragraph{Método (§5.7).} Variables:
$\Delta\phi_0\in\{60^{\circ},\dots,165^{\circ}\}$, barrido lineal
$\Delta\phi(t)=\Delta\phi_0+A(2t/T-1)$ con $A\in\{0,15^{\circ},30^{\circ},45^{\circ}\}$ y razón de radio
$\kappa\in\{0.85,1,1.15\}$, siempre con $z_{Rx}=z_{Tx}$: 96 candidatos.
Ninguna métrica contiene la potencia del eco. Se evalúan por separado
$Q_{img}=\mathrm{ISLR}_{mono}-\mathrm{ISLR}_{fus}$, $A_{PSF}=10\log_{10}(V_{fus}/V_{mono})$,
$\Delta\mathrm{SCR}=\mathrm{SCR}_{fus}-\mathrm{SCR}_{mono}$ (S1, $s=2$\,cm, TM, 2
realizaciones) y $C_{voo}=(v^{Rx}_{max}/v_{Tx}-1)_++|L_{Rx}/L_{Tx}-1|$
(inviable si la separación es $<15$\,m o $v^{Rx}_{max}>1.5\,v_{Tx}$).
Ponderación espectral Hann. El costo ponderado, con métricas normalizadas al
intervalo $[0,1]$ de los candidatos viables, es
\begin{equation}
  J = w_1\,\hat Q_{img} + w_2\,\widehat{\mathrm{SCR}} - w_3\,\hat A_{PSF} - w_4\,\hat C_{voo},
  \qquad (w_1,w_2,w_3,w_4)=(1,\,0.5,\,1,\,1).
\end{equation}
CSV: \texttt{expF\_candidatos.csv} (todas las métricas por candidato),
\texttt{expF\_eleccion\_pesos.csv}, \texttt{expF\_robustez.csv}, \texttt{expF2\_fase.csv}.

\begin{figure}[H]
\centering
\includegraphics[width=\textwidth]{F_pareto.pdf}
\caption{Frontera de Pareto (círculos rojos, 38 de 96 candidatos) en
las proyecciones $Q_{img}$--$C_{voo}$, $Q_{img}$--$\Delta$SCR y $Q_{img}$--$A_{PSF}$;
color $=\Delta\phi_0$; estrella: máximo de $J$.}
\label{fig:F_pareto}
\end{figure}

\begin{table}[H]
\centering\small
\caption{Elección según los pesos (Pareto-óptimas en todos los casos).}
\label{tab:F_pesos}
\resizebox{\textwidth}{!}{\begin{tabular}{@{}lcccc|ccc|rrrr@{}}
\toprule
Pesos & $w_1$ & $w_2$ & $w_3$ & $w_4$ & $\Delta\phi_0$ & $A$ & $\kappa$ & $Q_{img}$ & $\Delta$SCR & $A_{PSF}$ & $C_{voo}$ \\
\midrule
principal & 1 & 0.5 & 1 & 1 & 105° & 45° & 0.85 & 7.72 & 4.65 & 1.83 & 0.04 \\
sin SCR ni vuelo & 1 & 0 & 1 & 0 & 105° & 45° & 1.15 & 7.36 & 0.66 & 1.55 & 0.60 \\
más SCR & 1 & 1 & 0.5 & 1 & 105° & 45° & 0.85 & 7.72 & 4.65 & 1.83 & 0.04 \\
vuelo prioritario & 0.5 & 0.5 & 0.5 & 2 & 75° & 0° & 1.00 & 4.00 & 3.12 & 0.96 & 0.00 \\
\bottomrule
\end{tabular}}
\end{table}

\begin{table}[H]
\centering\small
\caption{Trayectoria propuesta frente a mono, a la mejor separación constante y a la región de Ref1
(misma ponderación Hann; métricas separadas, como exige §5.7).}
\label{tab:F_final}
\resizebox{\textwidth}{!}{\begin{tabular}{@{}lrrrrrrr@{}}
\toprule
 & PSLR & ISLR & $V_{-3}$ & $\delta_x$ & $\delta_z$ & Contr. & SCR \\
 & [dB] & [dB] & [l] & [m] & [m] & [dB] & [dB] \\
\midrule
Mono (Ref2) & $-7.73$ & 6.71 & 29.4 & 0.158 & 2.08 & 19.5 & 6.9 \\
Mono+bi, $\Delta\phi=105^{\circ}$ constante & $-10.40$ & 1.86 & 43.2 & 0.19 & 2.08 & 22.4 & 10.0 \\
\textbf{Mono+bi propuesta} ($60^{\circ}\to150^{\circ}$, $\kappa=0.85$) & $\mathbf{-10.10}$ & $\mathbf{-1.01}$ & 44.8 & 0.20 & 2.06 & \textbf{25.0} & 11.5 \\
Mono+bi, $\Delta\phi=165^{\circ}$ constante (región Ref1) & $-5.2$ & 7.5 & 61.7 & 0.23 & 2.08 & 14.5 & $-32.8$ \\
\bottomrule
\end{tabular}}
\end{table}

\begin{figure}[H]
\centering
\includegraphics[width=\textwidth]{F_trayectoria_final.pdf}
\caption{Trayectoria propuesta del Rx comparada con Ref1 y Ref2; a la
derecha, $\Delta\phi(t)$ con el Tx a 10\,m/s (190\,s de adquisición).}
\label{fig:F_final}
\end{figure}

\paragraph{Resultados e interpretación.}
\begin{enumerate}[itemsep=3pt]
  \item \textbf{Trayectoria propuesta:} Rx a la misma altura que el Tx, radio
  $0.85\,\rho_{Tx}$, separación que crece linealmente de 60° a 150° durante las
  dos vueltas. Frente a la imagen mono: ISLR $-7.7$\,dB, PSLR $-2.4$\,dB,
  contraste $+5.5$\,dB, a costa de un lóbulo principal $+52$\,\% más
  voluminoso ($\delta_x$ 0.158$\to$0.20\,m). Frente a la mejor separación
  constante (105°), el barrido gana 2.9\,dB de ISLR: los anillos de distinto
  radio rellenan el disco de $k$.
  \item \textbf{Costo de vuelo bajo:} velocidad máxima del Rx 0.97 veces la del
  Tx, recorrido 0.96, aceleración máxima 0.77\,m/s$^2$ a 10\,m/s, separación
  mínima 128\,m. Las 96 trayectorias son viables con las restricciones impuestas.
  \item \textbf{El radio no aporta resolución vertical (resultado negativo).} A la
  misma altura, $\kappa=0.85$--$1.15$ cambia $\delta_z$ en $\approx1$\,\%
  (2.06--2.08\,m): la extensión $k_z$ la fija la banda y no la diferencia de
  elevación ($\Delta\theta\approx\pm3^{\circ}$). $\kappa$ sólo influye en $C_{voo}$ y
  marginalmente en $Q_{img}$.
  \item \textbf{SCR:} la ganancia de SCR de la fusión está entre $-0.5$ y $+5.6$\,dB para
  $\Delta\phi_0\le150^{\circ}$ y es del orden de la dispersión entre realizaciones
  ($\pm3$\,dB); sólo la región $\Delta\phi_0\ge165^{\circ}$ es claramente mala
  ($-33$ a $-40$\,dB, especular). El SCR discrimina la región Ref1, pero no
  ordena las demás.
  \item \textbf{Sensibilidad a los pesos:} con tres de cuatro juegos de pesos se
  elige el mismo barrido (105°, $A=45^{\circ}$); sólo si el vuelo domina se elige
  $\Delta\phi=75^{\circ}$ constante. La elección es estable.
\end{enumerate}

\section{Robustez, sincronización y calibración de fase}
\label{sec:sincro}

\begin{figure}[H]
\centering
\includegraphics[width=\textwidth]{F_robustez.pdf}
\caption{Robustez de la fusión con la trayectoria propuesta frente a errores del
canal bi \emph{sin calibrar}.}
\label{fig:F_rob}
\end{figure}

\begin{figure}[H]
\centering
\includegraphics[width=\textwidth]{F2_fase.pdf}
\caption{Fase residual constante entre canales: sin calibrar (naranja),
calibrada sobre el lóbulo principal del dispersor más brillante (verde),
fusión incoherente $\tfrac12(|I_m|^2+|I_b|^2)$ (discontinua) y mono (azul).}
\label{fig:F2}
\end{figure}

\begin{table}[H]
\centering\small
\caption{Tolerancias del canal bi (fusión propuesta; nominal: PSLR $-10.1$,
ISLR $-1.0$\,dB; mono: $-7.7$ / $6.7$\,dB).}
\label{tab:tolerancias}
\begin{tabular}{@{}L{4.7cm}L{4.0cm}L{5.6cm}@{}}
\toprule
Perturbación & Tolerada & Observación \\
\midrule
Ruido de posición del Rx (blanco, por pulso) & $\le10$\,cm ($\lambda_0/7$) & ISLR $-0.2$\,dB a 10\,cm; se promedia en la apertura \\
Deriva de posición (paseo aleatorio, rms) & $\le5$\,cm & 20\,cm destruye la ganancia (PSLR $-2.0$\,dB) \\
Ruido de fase del oscilador (por pulso) & $\le45^{\circ}$ rms & sólo pérdida de amplitud \\
Offset de sincronismo constante & \textbf{calibrable} & 1\,ns = 153° de fase: sin calibrar anula la fusión;
                                                       calibrado, PSLR $-9.7$, ISLR $-1.0$\,dB \\
Deriva de sincronismo (lineal) & $<1$\,ns en la adquisición & 1\,ns: PSLR $-6.5$\,dB; 5\,ns: $-2.1$\,dB \\
Error del índice del suelo ($\pm10$\,\%) & sí & afecta igual a mono y fusión (localización 0.17\,m) \\
\bottomrule
\end{tabular}
\end{table}

\paragraph{Requisitos de sincronización y calibración del Rx remoto (§11).}
\begin{itemize}[itemsep=2pt]
  \item La fusión coherente exige conocer \textbf{una fase constante} entre
  canales con error $\lesssim30^{\circ}$ (Fig.~\ref{fig:F2}: a 60° se pierde ya
  2.1\,dB de ISLR; a 150° la fusión es peor que la mono). Un offset de reloj
  constante es exactamente eso, más un desplazamiento en rango despreciable
  ($c\,\Delta t/2n$). Se estima sin ambigüedad sobre el lóbulo principal de un
  reflector fuerte (o con la señal directa Tx$\to$Rx, cuya geometría se conoce por
  GNSS-RTK) y la fusión calibrada recupera exactamente la nominal. Estimarla sobre
  toda la imagen falla ($\approx180^{\circ}$ de error) porque dominan los anillos en
  contrafase.
  \item Las \textbf{derivas} (de reloj o de posición) no son calibrables con una
  constante: la deriva de reloj debe mantenerse $\ll1$\,ns en los 190\,s de
  adquisición ($\lesssim5\times10^{-12}$ de estabilidad relativa, nivel
  GPSDO/OCXO disciplinado o enlace de sincronización bidireccional), y la deriva de posición por debajo de $\approx5$\,cm rms
  (RTK). El ruido blanco de posición y de fase se tolera bien.
  \item La \textbf{fusión incoherente} (sumar intensidades) no necesita fase pero
  rinde 2.0\,dB menos de ISLR y 4.4\,dB menos de contraste que la coherente
  calibrada; su PSLR es similar ($-10.7$\,dB). Es un respaldo si no hay calibración
  fiable, no un sustituto.
\end{itemize}

\conclexp{F}{La trayectoria recomendada (Rx a la misma altura, radio $0.85\rho_{Tx}$ y separación
barrida de 60° a 150°) baja el ISLR 7.7\,dB y sube el contraste 5.5\,dB respecto de la mono con
un costo de vuelo casi nulo, siempre que la fase entre canales se calibre: un offset de reloj de
1\,ns sin calibrar anula la fusión.}

% ══════════════════════════════════════════════════════════════════════════
\chapter{Respuestas a las preguntas de investigación}
\label{cap:respuestas}

\begin{longtable}{@{}L{4.3cm}L{10.2cm}@{}}
\caption{Respuestas (§4) y criterio de éxito (§12), con la evidencia que las sustenta.}
\label{tab:respuestas}\\
\toprule
\textbf{Pregunta} & \textbf{Respuesta} \\
\midrule
\endfirsthead
\toprule
\textbf{Pregunta} & \textbf{Respuesta} \\
\midrule
\endhead
1. ¿Cómo altera $\Delta\phi$ el soporte en $k$ y la PSF? &
A la misma altura y radio, el canal bi ocupa un anillo horizontal de radio
$|\cos(\Delta\phi/2)|$ veces el mono, con idéntica extensión en $k_z$
(Figs.~\ref{fig:B_k}, \ref{fig:A_k}). La PSF bi se ensancha como
$\delta_{xy}/|\cos(\Delta\phi/2)|$ (error $<1$\,\% hasta 150°) y su resolución
vertical no cambia. En Ref1 el anillo colapsa en el origen y no hay
enfoque horizontal ($\delta_x=17$\,m). \\
2. ¿Mejor imagen bi aislada? &
La de menor separación viable: toda separación degrada el bi aislado; no hay
óptimo intermedio. Con la separación mínima de seguridad el bi aislado
coincide prácticamente con la mono y no aporta información. En SCR, en cambio, el bi aislado
en TM a 135°--140° es la mejor imagen con blancos reales: 8 a 20\,dB sobre la mono (Exp.~C2),
con una resolución horizontal 2.6 veces peor. \\
3. ¿Qué separación añade más información a la mono? &
$\Delta\phi\approx90^{\circ}$--$120^{\circ}$ constante (rango efectivo de la matriz de
adquisición 28.4$\to$33.8; área ocupada en $(k_x,k_y)$ $+52$\,\% a 90° y $+62$\,\% a 75°); mejor aún un barrido
$60^{\circ}\to150^{\circ}$ (Exp.~F). Fuera de $[60^{\circ},150^{\circ}]$ la complementariedad cae. \\
4. ¿Mejora la fusión coherente resolución, volumen, PSLR, ISLR, contraste, localización? &
\textbf{Mejora lóbulos y contraste, no la resolución.} Punto isótropo,
fusión N: PSLR $-2.0$ (90°) a $-5.2$\,dB (120°), ISLR $-5.3$\,dB (90°), contraste
$+4.0$\,dB (90°); con el barrido del Exp.~F: ISLR $-7.7$\,dB, contraste $+5.5$\,dB.
\textbf{Empeora} la resolución $-3$\,dB horizontal ($+16$ a $+30$\,\%) y el volumen
del lóbulo ($+35$ a $+58$\,\%); la vertical no cambia. La localización del punto ya es
perfecta en mono (error $<2$\,cm) y no cambia. La separabilidad de blancos próximos
empeora (73\,\% $\to$ 27\,\% a 0.30\,m). Todo es geométrico: fusN $\equiv$ fus. \\
5. ¿Se mantiene con RCS anisótropa y orientaciones? &
\textbf{Sí, pero sólo con un Rx de polarización alineada y a 45°--90°.}
Con Rx TM co-polar, a 90° el canal bi de un objeto real cae $\approx20$\,dB en
potencia media (desajuste geométrico de polarización en el suelo) y la fusión de igual
peso sólo deja $+0.3$ a $+0.6$\,dB de contraste a costa de 5--6\,dB de SNR. Con el Rx
alineado (doble polarización), sobre 200--1000 orientaciones SO(3) por adquisición, la
fusión a 90° mejora el contraste $+1.8$ a $+2.7$\,dB con probabilidad 90--98\,\% en cubo,
caja, cilindro y tubo, con $-1$ a $-2$\,dB de SNR (N). En orientación nominal, los objetos
con cara horizontal hacen que el especular bi domine a 135°--180° y la fusión empeora.
En escenas con varios objetos, la fusión a 135° reduce el enmascaramiento 10\,dB.
\\
6. ¿El pseudo-Brewster reduce el retorno superficial de forma robusta? &
\textbf{Sólo parcialmente y sólo donde no interesa.} En S1 el mínimo de reflectancia
TM es ancho pero poco profundo ($-34.5$\,dB; pierde 0.7/2.4/6.0/13\,dB con 0.5/1/2/5°
de error; roll 5° irrelevante): es robusto a errores, pero sólo suprime 17\,dB respecto de
la incidencia base. Esa supresión actúa sobre la especular coherente, que sólo llega al
Rx en la geometría Ref1: allí el SCR mejora 14--16\,dB pero queda en $\approx-41$\,dB (el blanco
sigue enmascarado). En mono y en las separaciones intermedias el clutter dominante es
difuso y el SCR no depende de Brewster (mismas tendencias en TM y TE). En el suelo sin
pérdidas el cero es exacto pero desaparece con 0.5°. En suelos húmedos el mínimo es más
somero ($-30.7$ y $-26.5$\,dB) y la atenuación del blanco domina. El efecto útil es
otro: el nulo SPM ``Brewster biestático'' en VV ($\Delta\phi\approx135^{\circ}$--$140^{\circ}$), que con
esfera y cubo en orientación aleatoria da al canal bi aislado 8 a 20\,dB más de SCR que la mono
(Exp.~C2). El nulo HH en TE a 90° no sirve: también desacopla el eco del blanco real. Ojo: el
nulo VV es de primer orden; el segundo orden lo rellena.
\\
7. ¿Costo en imagen de imponer pseudo-Brewster? &
\textbf{Pequeño:} la órbita de incidencia constante en $\theta_{pB}$ mejora la
resolución horizontal un 11\,\%, empeora la vertical un 12.5\,\%, sube el ISLR 1.2--1.8\,dB,
baja el contraste 0.5--1.2\,dB y baja el SCR 0.6--2.2\,dB. Además, la órbita de
incidencia constante (F2) da \emph{menos} SCR que la hélice (F1) en 54 de 60 comparaciones
($-3.5$\,dB de media; Tabla~\ref{tab:C_f1f2}). El único beneficio aparece en Ref1, inservible
para imagen. No compensa como restricción de trayectoria (Tabla~\ref{tab:C_concl}).
\\
\midrule
\multicolumn{2}{@{}l}{\textbf{Criterio de éxito (§12)}} \\
¿El canal bi añade información a la mono? &
Sí, en forma de lóbulos secundarios más bajos, mayor contraste, mejor
condicionamiento y menor enmascaramiento entre objetos de RCS distinta; no en
resolución. Sólo con calibración de fase entre canales. \\
¿Qué posición/trayectoria lo maximiza? &
Rx a la misma altura, $\kappa\approx0.85$, $\Delta\phi(t)$ barrido $60^{\circ}\to150^{\circ}$
(o, si el vuelo manda, $\Delta\phi=75^{\circ}$--$90^{\circ}$ constante), con Rx de doble
polarización alineado al campo del Tx. Para suprimir clutter, en cambio, $\Delta\phi\approx140^{\circ}$
constante en TM, con el canal bi aislado o una fusión ponderada. Nunca Ref1. \\
¿Dónde está la mejora? & Lóbulos (PSLR/ISLR) y contraste con la fusión; supresión de
clutter con el canal bi aislado en TM a 135°--140° (ver pregunta 6). \\
¿Se mantiene para blancos, orientaciones y suelos? &
Para orientaciones aleatorias, sí en contraste (con Rx alineado); no en resolución.
Para suelos, la fusión es geométrica e independiente del suelo (B, C: mismas tendencias en
S0--S1), pero en suelos húmedos (S2, S3) el SCR del blanco es negativo en todas las
geometrías.
\\
¿Vale la pena imponer pseudo-Brewster? &
No como restricción: la ganancia de SCR se limita a Ref1 y desaparece en suelos húmedos o
con 0.5° de error en el caso ideal, y la órbita que la impone es peor que la hélice en SCR y
resolución vertical. El SCR sube con la incidencia y con el Rx a 135°--140° en TM usado como canal aislado,
no con Brewster.
\\
\bottomrule
\end{longtable}

% ══════════════════════════════════════════════════════════════════════════
\chapter{Resultados negativos}
\label{cap:negativos}

Lista explícita (§10). Cada punto es un resultado reproducible del CSV indicado.
\begin{enumerate}[itemsep=3pt]
  \item \textbf{Ref1 (diametral) no sirve para imagen} y su fusión con la mono es
  peor que la mono sola (Exp.~A, B): PSLR $-7.9\to-3.8$\,dB, ISLR $+5.6$\,dB.
  \item \textbf{La fusión no mejora la resolución $-3$\,dB} en ninguna separación
  (Exp.~B, F): la ensancha. La anisotropía 3D baja sólo porque el lóbulo
  horizontal se ensancha.
  \item \textbf{El Rx a la misma altura no aporta resolución vertical}, ni con
  $\Delta\phi$ ni con el radio $\kappa=0.85$--$1.15$ ($\delta_z$ cambia $\le1$\,\%).
  \item \textbf{La separabilidad de blancos próximos empeora} con la fusión
  (Exp.~D): a 0.30\,m, 73\,\% (mono) frente a 27\,\% (90°).
  \item \textbf{La fusión de igual peso diluye la SNR} cuando el canal bi es débil:
  objetos con caras planas a 90° pierden 5--7\,dB de SNR (con el mismo presupuesto N).
  \item \textbf{Más eco biestático no implica mejor imagen}: a 135°--180° el canal bi
  de cubo/caja/cilindro es $+6$ a $+30$\,dB más fuerte y la fusión hereda su lóbulo
  ancho ($\Delta V$ hasta $+28$\,dB).
  \item \textbf{La fusión coherente sin calibrar puede ser peor que la mono}: una fase
  residual de 150°--180° (p.\,ej.\ 1\,ns de offset de reloj) la anula.
  \item \textbf{El SCR no ordena las separaciones intermedias} del Exp.~F (diferencias
  dentro de $\pm3$\,dB de dispersión entre realizaciones).
  \item \textbf{El pseudo-Brewster no rescata Ref1}: SCR $\approx-41$\,dB en S1 aun en el
  mínimo; sólo en el suelo ideal sin pérdidas Ref1 llega a $+61$\,dB, y lo pierde con $\pm2^{\circ}$.
  \item \textbf{El pseudo-Brewster no cambia el SCR de mono ni de las separaciones intermedias}
  (Exp.~C): sus variaciones con la incidencia son iguales en TE y en TM.
  \item \textbf{En suelos húmedos (S2, S3) ninguna geometría alcanza SCR $>0$} para un blanco
  de $1\,\mathrm{m}^2$ a 2\,m.
  \item \textbf{La órbita de altura constante acumula coherentemente los lóbulos lejanos en rango
  del reflejo de nadir} sobre su eje si no se pondera el espectro (Exp.~C, rectangular).
  \item \textbf{Un Rx TM co-polar anula la complementariedad a $\Delta\phi\approx90^{\circ}$}
  para objetos reales: $-20$\,dB de potencia media del canal bi (Exp.~D y E); el punto
  isótropo escalar no lo revela.
  \item \textbf{El nulo HH en TE a 90° no mejora el SCR de blancos reales} (Exp.~C2): el mismo
  desacople que anula el clutter anula el eco; la esfera pierde 7.1\,dB de SCR frente a la mono.
  El punto escalar, que daba $+17.8$\,dB, no sirve para evaluar efectos de polarización.
  \item \textbf{La fusión de igual peso desperdicia un canal con poco clutter} (Exp.~C2): a 140°
  en TM queda 10 a 15\,dB por debajo del canal bi aislado.
  \item \textbf{Con orientación aleatoria la fusión no mejora PSLR, resolución ni
  localización} (medianas $\approx0$; Exp.~E).
\end{enumerate}

% ══════════════════════════════════════════════════════════════════════════
\chapter{Limitaciones}
\label{cap:limitaciones}

\begin{itemize}[itemsep=2pt]
  \item Blancos PEC por PO (sin difracción de aristas, ondas de superficie ni
  rebotes múltiples); el tubo ($ka\approx2.7$) está en el límite de validez de PO.
  Co-polar únicamente (sin polarización cruzada del blanco ni del SPM fuera del plano).
  \item Superficie: SPM de primer orden ($k_0s\le0.36$) más especular coherente;
  sin términos de segundo orden, que rellenan los nulos del SPM (el nulo VV
  ``biestático'', del que depende la ventaja del bi aislado a 135°--140°, es por tanto una cota
  optimista). Facetas
  en $\pm12$\,m: el clutter más lejano se omite.
  \item Suelo homogéneo, $\varepsilon$ constante en la banda; geometría de rayos con
  $\mathrm{Re}(n)$; sin clutter volumétrico ni capas.
  \item Ganancia de antena gaussiana apuntada al centro; sin acoplamiento directo
  Tx--Rx ni su supresión; ruido térmico sin figura de ruido.
  \item SCR en imagen estimado con 2--3 realizaciones de la superficie congelada
  (dispersión $\approx\pm3$\,dB) en la vecindad del blanco.
  \item La hélice del Tx es la óptima del trabajo anterior y es fija; con dos vueltas
  completas la PSF es casi invariante en la escena, por lo que las conclusiones sobre
  $\Delta\phi$ podrían cambiar con aperturas parciales.
\end{itemize}

% ══════════════════════════════════════════════════════════════════════════
\chapter{Conclusiones y propuesta para la reunión}
\label{cap:conclusiones}

\section{Conclusiones}
\begin{enumerate}[itemsep=3pt]
  \item \textbf{El canal biestático sí añade información a la imagen monoestática}, pero no
  donde se buscaba antes: no en potencia (confirmado otra vez: a 180° la potencia media bi
  es igual a la mono), ni en resolución $-3$\,dB (la fusión la ensancha), ni en resolución
  vertical (a la misma altura $k_z$ no cambia). La aporta en \textbf{lóbulos secundarios,
  contraste, condicionamiento y enmascaramiento}: con un punto, la fusión a 90°--120° baja
  el ISLR 5.3\,dB y el PSLR hasta 5.2\,dB; con la trayectoria barrida $60^{\circ}\to150^{\circ}$, el ISLR
  7.7\,dB y el contraste $+5.5$\,dB; con objetos reales y orientación aleatoria, el contraste
  $+1.8$ a $+2.7$\,dB (Rx alineado); en una escena con tres objetos, el enmascaramiento
  10\,dB.
  \item \textbf{Todo el beneficio es geométrico} (fusión con el mismo presupuesto de muestras
  idéntica a la de 2N) y \textbf{exige coherencia}: una fase constante entre canales debe
  calibrarse ($\lesssim30^{\circ}$); un offset de reloj de 1\,ns sin calibrar anula la fusión. La
  fusión incoherente (sumar potencias) rinde 2\,dB menos de ISLR y 4.4\,dB menos de contraste.
  \item \textbf{La región útil es $\Delta\phi\approx60^{\circ}$--$150^{\circ}$, con Rx de polarización
  alineada.} Ref1 (180°) es inservible y empeora la fusión; por debajo de 45° el canal bi es
  redundante. Un Rx TM co-polar pierde $\approx20$\,dB a 90° con objetos reales.
  \item \textbf{El pseudo-Brewster no justifica restringir la trayectoria.} En suelo con
  pérdidas el mínimo es ancho ($\pm1^{\circ}$ cuesta 2.4\,dB) pero somero ($-34.5$\,dB en S1) y
  sólo actúa sobre la especular coherente, que sólo llega en Ref1 (SCR $\approx-41$\,dB en el
  mínimo). En mono y en separaciones intermedias el clutter es difuso y el SCR depende de la
  rugosidad y la incidencia por igual en TM y TE. En el suelo sin pérdidas el mínimo es
  perfecto pero desaparece con 0.5°. El costo en imagen de la restricción es pequeño
  ($\pm12$\,\% en resoluciones).
  \item \textbf{Un nulo biestático del clutter es más útil que Brewster, pero sólo en TM.} El
  nulo SPM de VV en $\Delta\phi\approx135^{\circ}$--$140^{\circ}$ sobrevive con blancos reales: el canal
  bi aislado mejora el SCR 8 a 20\,dB sobre la mono con esfera y cubo en orientación aleatoria,
  a costa de una resolución horizontal 2.6 veces peor. El nulo HH en TE a 90°, que con el punto
  escalar parecía el mejor resultado, es un artefacto: también desacopla el eco del blanco. El
  nulo VV es de primer orden y debe verificarse con modelos de orden superior y medidas.
  \item \textbf{El canal bi tiene dos usos distintos.} Para imagen (lóbulos y contraste), fusión
  coherente con $\Delta\phi$ barrido de 60° a 150°. Para detección limitada por clutter, canal bi
  aislado en TM a $\approx140^{\circ}$. La fusión de igual peso sirve para lo primero, pero
  desperdicia lo segundo: hay que estudiar una fusión ponderada o una reconstrucción conjunta.
  \item \textbf{Resultados negativos} (Cap.~\ref{cap:negativos}): la contribución fuerte no es
  ``el biestático es mejor'', sino \emph{cuándo} (separación intermedia, Rx alineado,
  calibración de fase), \emph{dónde} (lóbulos, contraste, enmascaramiento) y \emph{por qué}
  (relleno del disco horizontal del dominio $k$).
\end{enumerate}

\section{Propuesta para la reunión del 17/10}
\begin{itemize}[itemsep=2pt]
  \item \textbf{Función de costo.} Adoptar
  $J=w_1\hat Q_{img}+w_2\widehat{\mathrm{SCR}}-w_3\hat A_{PSF}-w_4\hat C_{voo}$ con
  $Q_{img}=$ reducción de ISLR (o contraste) de la fusión, $A_{PSF}$ = crecimiento del
  volumen $-3$\,dB, $C_{voo}$ cinemático, sin término de potencia; pesos (1, 0.5, 1, 1). La
  elección es estable frente a los pesos. Proponemos sustituir la PSF del punto por el
  contraste medio del Monte Carlo SO(3) con Rx alineado como $Q_{img}$ en la siguiente
  iteración.
  \item \textbf{Trayectoria propuesta.} Para imagen: Rx a la misma altura, radio
  $0.85\rho_{Tx}$, $\Delta\phi(t)$ de 60° a 150° en las dos vueltas, receptor de doble
  polarización y calibración de fase sobre la señal directa o un reflector. Para detección con
  clutter fuerte: $\Delta\phi\approx140^{\circ}$ constante en TM. Falta validar la primera con
  objetos anisótropos, orientaciones aleatorias y clutter, porque se eligió con el punto.
  \item \textbf{Próximo artículo.} ``Physics-aware receiver-trajectory design for complementary
  monostatic--bistatic subsurface UAV-SAR imaging'': contribución en la fusión coherente
  (relleno del dominio $k$), el nulo biestático VV como criterio de ubicación del Rx frente al
  clutter, la regla de polarización del Rx y los requisitos de sincronización; el pseudo-Brewster pasaría a ser un análisis de sensibilidad
  (resultado negativo documentado) y no un criterio de diseño.
  \item \textbf{Preguntas abiertas.} (i) ¿Aceptamos perder 16--30\,\% de resolución $-3$\,dB a
  cambio de 5--8\,dB de ISLR? (ii) ¿Qué apertura parcial (menos de dos vueltas) es realista?
  Con aperturas parciales el dominio $k$ deja de ser anular y el óptimo de $\Delta\phi$ puede
  cambiar. (iii) ¿Disponemos de Rx de doble polarización en el RD350? (iv) Validar el nulo
  VV con un modelo IEM/AIEM, polarización cruzada y medidas. (v) ¿Qué combinación de canales
  conviene: fusión ponderada, reconstrucción conjunta o incoherente?
\end{itemize}


\appendix
% ══════════════════════════════════════════════════════════════════════════
\chapter{Resolución mono y biestática en forma cerrada}
\label{ap:resolucion}

Este apéndice da las expresiones cerradas de $\delta_x$, $\delta_y$ y $\delta_z$ para un
blanco enterrado en una posición arbitraria (dentro o fuera del eje de la órbita) y
para cualquier geometría Tx--Rx. Las compara con las resoluciones medidas en la imagen.
En todo el relatorio, $\delta_x$, $\delta_y$, $\delta_z$ denotan el ancho a $-3$\,dB de
$|\chi|$ (la PSF en amplitud) en cortes por el pico. Equivale al FWHM en potencia; en los
CSV las columnas se llaman \texttt{fwhm\_x}, \texttt{fwhm\_y}, \texttt{fwhm\_z}. Código:
\texttt{bisar/resolucion.py}; verificación: \texttt{t\_resolucion\_analitica.py}.

\section{Modelo}

Sea el blanco $\vect P=(x_P,y_P,z_P)$ a profundidad $d=-z_P$ en un suelo de índice
$n=\mathrm{Re}\sqrt{\varepsilon_r}$, y sean el Tx y el Rx en $\vect p_i=(x_i,y_i,h_i)$,
$i\in\{T,R\}$. Para cada pulso $m$ y frecuencia $f$, la fase del eco tras la
back-projection en $\vect P+\Delta\vect r$ es $\vect K_{m,f}\cdot\Delta\vect r$, con el
vector de onda biestático
\begin{equation}
  \vect K_{m,f}=\frac{2\pi f}{c}\,n\,\big(\hat{\vect u}_{T,m}+\hat{\vect u}_{R,m}\big),
  \qquad
  \chi(\Delta\vect r)=\sum_{m}\sum_{f} e^{\,j\vect K_{m,f}\cdot\Delta\vect r},
  \label{eq:A_psf}
\end{equation}
donde $\hat{\vect u}_i$ es el vector unitario, dentro del suelo, desde el blanco hacia el
punto de refracción de la plataforma $i$. Se toman pesos uniformes (ventana espectral
rectangular, como en el Exp.~B).

\paragraph{Geometría vista desde el blanco.} Para cada pulso,
\begin{gather}
  D_i=\sqrt{(x_i-x_P)^2+(y_i-y_P)^2},\qquad
  \phi'_i=\operatorname{atan2}(y_i-y_P,\;x_i-x_P),\label{eq:A_geo}\\
  D_i=h_i\tan\psi_i+d\tan\theta_i,\qquad \sin\psi_i=n\sin\theta_i .\label{eq:A_snell}
\end{gather}
$\psi_i$ es la incidencia en el aire y $\theta_i$ la refractada. La ecuación de Snell
(\ref{eq:A_snell}) se resuelve en forma explícita con un paso de corrección:
\begin{equation}
  \psi_i^{(0)}=\arctan\frac{D_i}{h_i},\qquad
  \psi_i=\arctan\!\left(\frac{D_i-d\,\sin\psi_i^{(0)}\big/\sqrt{n^2-\sin^2\psi_i^{(0)}}}{h_i}\right).
  \label{eq:A_psi}
\end{equation}
Con $d=2$\,m y $h\approx94$--$131$\,m, el resultado coincide con la solución exacta en
cuatro cifras. Entonces
$\hat{\vect u}_i=(\sin\theta_i\cos\phi'_i,\ \sin\theta_i\sin\phi'_i,\ \cos\theta_i)$.

\section{Resolución horizontal}

Como $n\sin\theta_i=\sin\psi_i$, la componente horizontal de (\ref{eq:A_psf}) no depende
del suelo:
\begin{equation}
  \vect K_h=k_0\big[\sin\psi_T\,\hat{\vect e}(\phi'_T)+\sin\psi_R\,\hat{\vect e}(\phi'_R)\big],
  \qquad k_0=\frac{2\pi f}{c},
\end{equation}
\begin{equation}
  |\vect K_h|=2k_0\,S,\qquad
  S=\tfrac12\sqrt{\sin^2\psi_T+\sin^2\psi_R+2\sin\psi_T\sin\psi_R\cos(\phi'_R-\phi'_T)} .
  \label{eq:A_S}
\end{equation}
Para el monoestático $S=\sin\psi_T$. Con $\psi_T=\psi_R=\psi$, $S=\sin\psi\,|\cos(\Delta\phi'/2)|$.

La apertura recorre $360^{\circ}$ alrededor del blanco, de modo que el soporte horizontal es
un anillo de radio $R_c=2k_0\langle S\rangle$, y la PSF horizontal es
$\chi(\rho)\propto J_0(R_c\rho)$. El ancho a $-3$\,dB sale de $J_0(x_{J_0})=1/\sqrt2$, con
$x_{J_0}=1.1264$:
\begin{equation}
  \boxed{\;\delta_x=\delta_y=\frac{2x_{J_0}}{R_c}=\frac{x_{J_0}\,\lambda_0}{2\pi\,\langle S\rangle}
  =\frac{0.1793\,\lambda_0}{\langle S\rangle}\;}
  \label{eq:A_dxy}
\end{equation}
donde $\langle\cdot\rangle$ es la media sobre los pulsos de la apertura y $\lambda_0=c/f_c$.
Nota: \texttt{explicacion\_resolucion.tex} usa la constante 1.20; como
$J_0(1.20)=0.67$ ($-3.5$\,dB), esa constante sobreestima $\delta_{xy}$ en un 6.5\,\%.

\section{Resolución vertical}

La componente vertical no depende del acimut:
\begin{equation}
  K_z=\frac{2\pi f}{c}\,q,\qquad
  q=n(\cos\theta_T+\cos\theta_R)=\sqrt{n^2-\sin^2\psi_T}+\sqrt{n^2-\sin^2\psi_R}.
  \label{eq:A_q}
\end{equation}
El soporte en $k_z$ es el producto de dos dispersiones: la banda ($f\in f_c\pm B/2$) y la
variación de $q$ a lo largo de la apertura. Para una distribución casi uniforme de cada una,
la PSF vertical es el producto de dos funciones sinc. Igualando la curvatura del lóbulo
principal, el ancho a $-3$\,dB es el de una sola sinc con banda equivalente
$\sqrt{(B\bar q)^2+12f_c^2\sigma_q^2}$:
\begin{equation}
  \boxed{\;\delta_z=\frac{c_w\,c}{\sqrt{(B\,\bar q)^2+12\,f_c^2\,\sigma_q^2}}\;},\qquad
  \bar q=\langle q\rangle,\quad \sigma_q^2=\operatorname{var}(q),
  \label{eq:A_dz}
\end{equation}
con $c_w=0.886$ para espectro rectangular y $c_w=1.443$ para Hann. El primer término es la
resolución por banda, $c_w c/(B\bar q)$; el segundo, la apertura tomográfica en incidencia.

\section{Caso del Exp.~B en forma totalmente explícita}

Con el blanco en el eje y el Rx en la misma hélice desfasada $\Delta\phi$, se cumple
$\psi_R=\psi_T$ y $\phi'_R-\phi'_T=\Delta\phi$ en cada pulso. Si la incidencia se barre
linealmente entre $\psi_1$ y $\psi_2$, las medias de (\ref{eq:A_dxy}) y (\ref{eq:A_dz})
se integran en forma cerrada. Con $Q(\psi)=\sqrt{n^2-\sin^2\psi}$:
\begin{align}
  \delta_x=\delta_y&=\frac{x_{J_0}\,\lambda_0}{2\pi\,|\cos(\Delta\phi/2)|}\;
  \frac{\psi_2-\psi_1}{\cos\psi_1-\cos\psi_2},\label{eq:A_dxy_eje}\\
  \delta_z&=\frac{c_w\,c}{\sqrt{B^2\,[Q(\psi_1)+Q(\psi_2)]^2+4f_c^2\,[Q(\psi_1)-Q(\psi_2)]^2}} .
  \label{eq:A_dz_eje}
\end{align}
En (\ref{eq:A_dz_eje}) se usa que, para $q$ lineal en el tiempo, $12\sigma_q^2=\Delta q^2$ con
$\Delta q=2[Q(\psi_1)-Q(\psi_2)]$. La ecuación (\ref{eq:A_dz_eje}) no contiene $\Delta\phi$.

\paragraph{Sustitución numérica.} Hélice óptima vista desde $(0,0,-2)$\,m por
(\ref{eq:A_psi}): $\psi_1=46.11^{\circ}$, $\psi_2=60.37^{\circ}$. Con $f_c=425$\,MHz
($\lambda_0=0.7059$\,m), $B=50$\,MHz y $n=2.0$:
\begin{itemize}[itemsep=1pt]
  \item $\langle\sin\psi\rangle=(\cos\psi_1-\cos\psi_2)/(\psi_2-\psi_1)=0.7991$ y
  $x_{J_0}\lambda_0/2\pi=0.12655$\,m, de donde
  $\delta_{xy}^{mono}=0.158$\,m y $\delta_{xy}(90^{\circ})=0.158/\cos45^{\circ}=0.224$\,m;
  \item $Q(\psi_1)=1.8656$ y $Q(\psi_2)=1.8012$, de donde $B\bar q=1.833\times10^{8}$\,Hz y
  $f_c\Delta q=0.548\times10^{8}$\,Hz, con lo que
  $\delta_z=0.886\cdot3\times10^{8}/1.913\times10^{8}=1.389$\,m para todo $\Delta\phi$.
\end{itemize}
La imagen da 0.159, 0.223 y 1.387\,m (Tabla~\ref{tab:A_verif}).

\section{Blanco fuera del eje}

Las ecuaciones (\ref{eq:A_dxy}) y (\ref{eq:A_dz}) son válidas para cualquier $\vect P$; sólo
cambian las medias. A primer orden en $r_P/\rho$ (con $r_P=|\vect P_h|$ y $\rho$ el radio de
la órbita), para Tx y Rx en la misma hélice:
\begin{gather}
  \Delta\phi'(t)\simeq\Delta\phi-\frac{2r_P}{\rho}\,\sin\frac{\Delta\phi}{2}\,
  \cos\!\Big(\phi_T+\frac{\Delta\phi}{2}-\phi_P\Big),\label{eq:A_dphiap}\\
  q(t)\simeq q_0(t)-2\,\big|Q'(\psi)\big|\,\frac{\cos^2\psi}{h}\,r_P\,
  \cos\frac{\Delta\phi}{2}\,\cos\!\Big(\phi_T+\frac{\Delta\phi}{2}-\phi_P\Big),
  \qquad Q'(\psi)=-\frac{\sin\psi\cos\psi}{\sqrt{n^2-\sin^2\psi}} .\label{eq:A_qap}
\end{gather}
Hay dos consecuencias:
\begin{enumerate}[itemsep=1pt]
  \item La separación aparente oscila alrededor de $\Delta\phi$; en $(6,4,-2)$\,m y a $90^{\circ}$,
  la amplitud de (\ref{eq:A_dphiap}) es $3.9^{\circ}$ (medida: $86.1^{\circ}$--$94.3^{\circ}$). La
  oscilación tiene media nula, de modo que $\langle S\rangle$ y $\delta_{xy}$ no cambian a
  primer orden: fuera del eje vale la misma ley $\delta_{xy}^{mono}/|\cos(\Delta\phi/2)|$.
  \item La diferencia de incidencia entre el lado cercano y el lejano añade a $q$ una
  oscilación que aumenta $\sigma_q$ y mejora ligeramente $\delta_z$. En el canal bi esa
  oscilación se escala por $|\cos(\Delta\phi/2)|$, porque Tx y Rx están en lados distintos y
  sus variaciones se compensan. Por eso, fuera del eje, $\delta_z$ pasa de 1.378\,m (mono)
  a 1.386\,m a $90^{\circ}$ y vuelve al valor del eje (1.387\,m) a $180^{\circ}$.
\end{enumerate}

\section{¿Cuándo puede el biestático resolver mejor que la mono?}

\paragraph{Horizontal.} Por la desigualdad triangular en (\ref{eq:A_S}),
\begin{equation}
  S\le\tfrac12(\sin\psi_T+\sin\psi_R)
  \quad\Longrightarrow\quad
  \delta_{xy}^{bi}\ \ge\ \frac{x_{J_0}\,\lambda_0}{\pi\,\langle\sin\psi_T+\sin\psi_R\rangle} .
\end{equation}
Por tanto $\delta_{xy}^{bi}<\delta_{xy}^{mono,Tx}$ exige $\sin\psi_R>\sin\psi_T$: el Rx tiene
que ver el blanco a mayor incidencia (más bajo o más lejos) y con $\Delta\phi$ pequeño. Con
$\psi_R=\psi_T$ (Exp.~B), $S=\sin\psi\,|\cos(\Delta\phi/2)|\le\sin\psi$, y el bi es
\emph{siempre} igual (en $\Delta\phi=0$) o peor.

\paragraph{Vertical.} El término de banda de (\ref{eq:A_dz}) mejora si
$q=Q(\psi_T)+Q(\psi_R)>2Q(\psi_T)$, es decir, si $\psi_R<\psi_T$ (Rx más cerca del nadir).
Como $Q\le n$, la mejora está acotada por
\begin{equation}
  \frac{\delta_z^{mono}}{\delta_z^{bi}}\Big|_{banda}\le\frac{Q(\psi_T)+n}{2Q(\psi_T)}=1.046
  \qquad(\psi_T=53.2^{\circ},\ n=2):
\end{equation}
la refracción hace que $\cos\theta_t\ge0.92$ en el suelo, y la ganancia máxima es del 4.6\,\%. Con
$\psi_R=\psi_T$ (Exp.~B), $q$ es idéntico y $\delta_z^{bi}=\delta_z^{mono}$.

\paragraph{Compromiso.} Las condiciones son opuestas ($\psi_R>\psi_T$ para $x,y$;
$\psi_R<\psi_T$ para $z$) y ambas piden $\Delta\phi\approx0$, es decir, renunciar a la
diversidad angular que justifica el canal bi. Además, la mejor mono posible a la incidencia
del Rx resolvería todavía mejor. Las cuatro últimas filas de cada blanco en la
Tabla~\ref{tab:A_verif} lo confirman:
\begin{itemize}[itemsep=1pt]
  \item con $\kappa=1.3$ y con el Rx 60\,m más bajo, $\delta_{xy}$ mejora un 4\,\% y un 8\,\%
  pero $\delta_z$ empeora un 1.4\,\% y un 2.6\,\%;
  \item con el Rx casi en el nadir ($\kappa=0.2$), $\delta_z$ mejora un 1.7\,\% pero
  $\delta_{xy}$ empeora un 50\,\%.
\end{itemize}

\section{Verificación y límites de validez}

\begin{table}[H]
\centering\footnotesize\setlength{\tabcolsep}{3.6pt}
\caption{Verificación de las formas cerradas (\ref{eq:A_dxy}) y (\ref{eq:A_dz}) frente a las resoluciones medidas en la imagen (back-projection con refracción exacta, cortes finos de 2 y 5\,mm por el pico). $\varepsilon$: error relativo de la forma cerrada. CSV: \texttt{expB/expB\_resolucion\_analitica.csv}.}
\label{tab:A_verif}
\begin{tabular}{@{}llccccccccc@{}}
\toprule
Blanco & Rx & $\bar\psi_{Rx}$ & $\langle S\rangle$ & $\delta_x$ & $\delta_y$ & $\delta_{xy}^{\,\mathrm{cerr}}$ & $\varepsilon_x/\varepsilon_y$ [\%] & $\delta_z$ & $\delta_z^{\,\mathrm{cerr}}$ & $\varepsilon_z$ [\%] \\
 & & [°] & & [m] & [m] & [m] & & [m] & [m] & \\
\midrule
$(0,0,-2)$ & $\Delta\phi=0^{\circ}$ (mono) & 53.2 & 0.799 & 0.159 & 0.159 & 0.158 & -0.4/-0.5 & 1.387 & 1.389 & +0.1 \\
 & $\Delta\phi=45^{\circ}$ & 53.2 & 0.738 & 0.171 & 0.173 & 0.171 & +0.2/-1.1 & 1.387 & 1.389 & +0.1 \\
 & $\Delta\phi=90^{\circ}$ & 53.2 & 0.565 & 0.223 & 0.227 & 0.224 & +0.4/-1.3 & 1.387 & 1.389 & +0.1 \\
 & $\Delta\phi=135^{\circ}$ & 53.2 & 0.306 & 0.413 & 0.418 & 0.414 & +0.2/-1.0 & 1.387 & 1.389 & +0.1 \\
 & $\Delta\phi=165^{\circ}$ & 53.2 & 0.104 & 1.216 & 1.221 & 1.213 & -0.3/-0.6 & 1.387 & 1.389 & +0.1 \\
 & $\Delta\phi=0$, $\kappa=1.3$ & 60.1 & 0.832 & 0.153 & 0.153 & 0.152 & -0.5/-0.5 & 1.406 & 1.409 & +0.2 \\
 & $\Delta\phi=0$, $z_{Rx}=z_{Tx}-60$\,m & 70.6 & 0.870 & 0.146 & 0.146 & 0.146 & -0.6/-0.6 & 1.423 & 1.426 & +0.2 \\
 & $\Delta\phi=0$, $\kappa=0.2$ & 15.1 & 0.530 & 0.238 & 0.239 & 0.239 & +0.4/-0.1 & 1.364 & 1.371 & +0.5 \\
\midrule
$(6,4,-2)$ & $\Delta\phi=0^{\circ}$ (mono) & 53.2 & 0.799 & 0.159 & 0.160 & 0.158 & -0.2/-0.9 & 1.378 & 1.380 & +0.1 \\
 & $\Delta\phi=45^{\circ}$ & 53.2 & 0.738 & 0.171 & 0.174 & 0.171 & +0.4/-1.5 & 1.382 & 1.384 & +0.2 \\
 & $\Delta\phi=90^{\circ}$ & 53.2 & 0.565 & 0.223 & 0.227 & 0.224 & +0.4/-1.5 & 1.386 & 1.388 & +0.2 \\
 & $\Delta\phi=135^{\circ}$ & 53.2 & 0.306 & 0.414 & 0.418 & 0.414 & -0.0/-0.9 & 1.388 & 1.390 & +0.2 \\
 & $\Delta\phi=165^{\circ}$ & 53.2 & 0.104 & 1.213 & 1.219 & 1.213 & -0.0/-0.5 & 1.387 & 1.389 & +0.1 \\
 & $\Delta\phi=0$, $\kappa=1.3$ & 60.1 & 0.832 & 0.153 & 0.154 & 0.152 & -0.3/-0.9 & 1.399 & 1.402 & +0.2 \\
 & $\Delta\phi=0$, $z_{Rx}=z_{Tx}-60$\,m & 70.6 & 0.869 & 0.146 & 0.147 & 0.146 & -0.3/-1.1 & 1.416 & 1.420 & +0.3 \\
 & $\Delta\phi=0$, $\kappa=0.2$ & 15.4 & 0.531 & 0.237 & 0.240 & 0.238 & +0.5/-0.8 & 1.358 & 1.364 & +0.4 \\
\bottomrule
\end{tabular}
\end{table}


En los 34 casos con $\Delta\phi<180^{\circ}$ (dos blancos, Exp.~B completo y cinco geometrías
del Rx fuera del Exp.~B), las formas cerradas reproducen la imagen:
\begin{itemize}[itemsep=1pt]
  \item $|\varepsilon_x|\le0.6$\,\% y $|\varepsilon_y|\le1.6$\,\%; $\delta_y$ es hasta un 2\,\%
  mayor que $\delta_x$ porque la apertura de dos vueltas, con radio variable, no es
  perfectamente isótropa;
  \item $|\varepsilon_z|\le0.5$\,\%.
\end{itemize}
Las cifras del cuerpo del relatorio vienen de la rejilla 3D ($\Delta x=2.5$\,cm, con
interpolación) y difieren menos de un 1.5\,\% de estos cortes finos (p.\,ej. 0.157 frente a
0.159\,m en la mono).

Límites de validez:
\begin{enumerate}[itemsep=1pt]
  \item \textbf{Apertura completa.} (\ref{eq:A_dxy}) supone una apertura de $360^{\circ}$ y una
  dispersión de $|\vect K_h|$ pequeña frente a su media. Falla en $\Delta\phi\to180^{\circ}$: el
  anillo colapsa ($\langle S\rangle\to0$) y la anchura la fija la dispersión residual. Para el
  blanco en el eje, el corte largo da 17\,m (Exp.~A).
  \item \textbf{Pesos uniformes.} Se ignoran la ganancia de antena, $1/R^2$ y la transmitancia,
  que varían poco a lo largo de la hélice; con ventana de Hann se usa $c_w=1.443$.
  \item \textbf{Suelo.} Homogéneo, con $n=\mathrm{Re}\sqrt{\varepsilon_r}$: las pérdidas
  atenúan, pero no cambian la geometría de $\vect K$.
\end{enumerate}

% ══════════════════════════════════════════════════════════════════════════
\chapter{Reproducción, figuras y matriz de simulaciones}
\label{ap:repro}

\paragraph{Código.} \texttt{simulations/analise\_1610} (instrucciones en
\texttt{README.md}): \texttt{python run\_all.py} ejecuta A--F; además
\texttt{exp\_C\_figuras\_enlace.py}, \texttt{exp\_D2\_rejilla\_densa.py},
\texttt{exp\_F2\_calibracion\_fase.py}, \texttt{exp\_F\_fig\_robustez.py}. Un JSON
por experimento en \texttt{params/} (hereda de \texttt{common.json}); cada
corrida guarda \texttt{params\_resueltos\_expX.json}. Pruebas en \texttt{tests/}.
Todas las métricas en CSV en \texttt{simulations/io/analise\_1610/expX/}.

\begin{table}[H]
\centering\small
\caption{Conjunto mínimo de figuras (§7).}
\label{tab:figuras}
\begin{tabular}{@{}rL{7.4cm}L{5.4cm}@{}}
\toprule
\# & Figura pedida & En este relatorio \\
\midrule
1 & Geometría Tx, Rx, Ref1, Ref2, intermedia & Fig.~\ref{fig:geom} \\
2 & Soporte en $k$: Ref2, Ref1 y tres intermedias & Fig.~\ref{fig:B_k} (y \ref{fig:A_k}) \\
3 & PSF 3D y cortes $x,y,z$ mono, bi, mono+bi & Figs.~\ref{fig:A_cortes}, \ref{fig:B_img} \\
4 & Resoluciones $x,y,z$ vs.\ $\Delta\phi$ & Fig.~\ref{fig:B_res} \\
5 & Volumen y anisotropía vs.\ $\Delta\phi$ & Fig.~\ref{fig:B_vol} \\
6 & PSLR e ISLR vs.\ $\Delta\phi$ & Fig.~\ref{fig:B_lob} \\
7 & Reflectancia TM y TE vs.\ incidencia & Fig.~\ref{fig:C_refl} \\
8 & SCR vs.\ incidencia y error angular & Figs.~\ref{fig:C_scr}, \ref{fig:C_err} \\
9 & Imágenes mono, bi, fusión, misma escala & Figs.~\ref{fig:A_img}, \ref{fig:B_img}, \ref{fig:D_mip} \\
10 & Métricas de la fusión vs.\ mono por clase de blanco & Fig.~\ref{fig:D_fus}, \ref{fig:E_mc} \\
11 & Pareto imagen--SCR--vuelo & Fig.~\ref{fig:F_pareto} \\
12 & Trayectoria final vs.\ Ref1 y Ref2 & Fig.~\ref{fig:F_final} \\
\bottomrule
\end{tabular}
\end{table}

\begin{table}[H]
\centering\small
\caption{Matriz de simulaciones ejecutada (§8).}
\label{tab:matriz}
\begin{tabular}{@{}L{2.6cm}L{6.2cm}L{5.4cm}@{}}
\toprule
Eje & Valores pedidos & Ejecutados \\
\midrule
Canal & mono; bi; fusión & mono, bi, fus (2N), fusN (N), incoherente \\
Separación & 0°--180° en pasos de 15° & 13 valores (B); 8$\times$4$\times$3 (F) \\
Incidencia & fuera de Brewster, cerca, barrido local & 35, 45, 53.4, 75° y $\theta_{pB}+\{0,\pm0.5,\pm1,\pm2,\pm5\}$; dos familias \\
Suelo & sin pérdidas; $\varepsilon$ complejo; dos adicionales & S0, S1, S2, S3 \\
Superficie & lisa; rugosidad moderada & $s=0,1,2,4$\,cm, $\ell=20$\,cm \\
Polarización & TM y TE & TM, TE, roll 0--5° \\
Blanco & punto, esfera, cubo, caja, cilindro, tubería, rejilla 64 & los 7 + tubo y cilindro inclinados, rejilla densa, escena P2 \\
Orientación & nominal; inclinada; Monte Carlo SO(3) & nominal, inclinados, SO(3) por adquisición (200--1000) \\
Robustez & ángulo, trayectoria, fase, sincronismo, suelo & todas (Exp.~C y F) \\
\bottomrule
\end{tabular}
\end{table}


\end{document}
